\documentclass[10pt,a4paper]{article}

% ----------------- Paquetes -------------------%
\usepackage[T1]{fontenc}
\usepackage[utf8]{inputenc}
\usepackage[a4paper,margin=2.5cm]{geometry}
\usepackage{mathptmx}             
\usepackage{changepage} 
\usepackage{setspace}
\usepackage[spanish]{babel}
\usepackage[labelfont=bf,labelsep=space]{caption}
\usepackage{amssymb}
\usepackage{amsmath}
\usepackage{ebgaramond}
\usepackage{placeins}
\usepackage{etoolbox}
\usepackage{setspace}
\usepackage{parskip} 
\usepackage{tcolorbox}
\usepackage{needspace}
\usepackage{xparse}
\usepackage{ocgx2}
\usepackage{hyperref}
\usepackage{hypcap}
\usepackage{soul}\sethlcolor{yellow}% resaltado amarillo: \hl{texto} (Panel TCEE, ⌃H)



%%%%%%%%%%%%%%%%%%%%%%%%%%%%%%%%%%%%%%%%%%%%%%%%%%%%%%%%%%%%%%%%
% --- Fija primero tus valores globales "buenos" ---
\setstretch{1.2}                 % Interlineado global
\setlength{\parskip}{0.7\baselineskip}
\setlength{\parindent}{0pt}
\newlength{\GlobalParskip}
\setlength{\GlobalParskip}{\parskip}

\newcounter{eqnum}[subsection]
\renewcommand{\theeqnum}{\arabic{eqnum}}
\newcounter{alphaitem}
\newcounter{numitem}

%% ------------------- Recuadros----------------------%%
\newcommand{\puntosclave}[1]{%
  \begin{tcolorbox}[
    colframe=black,
    title={Puntos clave},
    before upper={\setlength{\parindent}{0.5cm}\setlength{\parskip}{0.5\baselineskip}}
  ]
  #1
  \end{tcolorbox}
}

\newcommand{\modificaciones}[1]{
  \begin{tcolorbox}[
    colback=white,
    colframe=red,
    title={Modificaciones a realizar},
    before upper={\setlength{\parindent}{0.5cm}\setlength{\parskip}{0.5\baselineskip}}
  ]
  #1
  \end{tcolorbox}
}

%% ------------------- Preguntas TEST----------------------%%
% Opciones: radio (single-choice)
\NewDocumentCommand{\OptRadio}{m m m}{%
  \noindent\CheckBox[radio,name=#1,value=#2,width=1.2ex,height=1.2ex]{}%
  \;\textbf{#2.}~#3\par
}

% Opciones: checkbox (multi-choice)
\NewDocumentCommand{\OptCheck}{m m}{%
  \noindent\CheckBox[name=#1,width=1.2ex,height=1.2ex,bordercolor={0 0 0}]{}%
  \;#2\par
}

% Pregunta single-choice (radio)
% Args: 1=id, 2=enunciado, 3=opciones, 4=solución+justificación, 5=imagen (o {}), 6=origen
\NewDocumentCommand{\PreguntaTestSingle}{m m m m m m}{%
  \par\medskip
  \noindent\textbf{#1.}~#2\par\smallskip

    % Imagen (si no es {})
    \IfBlankTF{#5}{}{%
      \begin{center}
        \includegraphics[width=0.75\linewidth]{#5}
      \end{center}
    }

    \begin{Form}
      #3
    \end{Form}

    \par\smallskip
    \switchocg{sol-#1}{\fbox{Mostrar / ocultar solución}}%
    \hfill{\footnotesize #6}\par\smallskip

    \begin{ocg}{Solución}{sol-#1}{0}
      \noindent #4
    \end{ocg}
  \par\medskip
}

% Pregunta multi-choice (checkboxes)
\NewDocumentCommand{\PreguntaTestMulti}{m m m m m m}{%
  \par\medskip
  \noindent\textbf{#1.}~#2\par\smallskip

    % Imagen (si no es {})
    \IfBlankTF{#5}{}{%
      \begin{center}
        \includegraphics[width=0.75\linewidth]{#5}
      \end{center}
    }
    \begin{Form}
      #3
    \end{Form}

    \par\smallskip
    \switchocg{sol-#1}{\fbox{Mostrar / ocultar solución}}%
    \hfill{\footnotesize #6}\par\smallskip

    \begin{ocg}{Solución}{sol-#1}{0}
      \noindent #4
    \end{ocg}
  \par\medskip
}

%% ------------------- Listas----------------------%%
    % --- Lista alfabética ---
    \newenvironment{la}
      {\begin{list}{\alph{alphaitem})}{%
          \usecounter{alphaitem}%
          \setlength{\leftmargin}{1cm}%
          \setlength{\parskip}{3pt}% 
          \setlength{\itemsep}{0pt}% ← espacio entre ítems
          \setlength{\parsep}{0pt}%  ← párrafos dentro de un ítem
      }}
      {\end{list}}
      
    % --- Lista numérica ---
    \newenvironment{lnum}
      {\begin{list}{\arabic{numitem}.}{%
          \usecounter{numitem}%
          \setlength{\leftmargin}{1cm}%
          \setlength{\parskip}{3pt}% 
          \setlength{\itemsep}{0pt}% ← espacio entre ítems
          \setlength{\parsep}{0pt}%  ← párrafos dentro de un ítem
      }}
      {\end{list}}

%% ------------------- Preguntas Test ----------------------%%
      \usepackage{xparse} % si ya lo tienes, no lo repitas

    \NewDocumentEnvironment{test}{m m o}{%
      \begin{tcolorbox}[
        colback=white,
        colframe=black!15,
        boxrule=0.6pt,
        arc=2mm,
        left=2mm,right=2mm,top=2mm,bottom=2mm
      ]
      \centering
      \includegraphics[width=\linewidth]{#1}
      \end{tcolorbox}
    
      \vspace{2mm}
    
      % Caja SOLO para texto (SÍ partible y mantiene márgenes al partir)
      \begin{tcolorbox}[
        enhanced,
        breakable,
        colback=black!3,
        colframe=black!25,
        boxrule=0.6pt,
        arc=2mm,
        left=2mm,right=2mm,top=1.5mm,bottom=1.5mm
      ]
      \textbf{Respuesta:} \textbf{#2}%
      \IfValueT{#3}{\par\smallskip \textbf{Justificación:} #3}
      \end{tcolorbox}
    }{}
      
% ------------------- Imágenes----------------------%
    \graphicspath{{Img/}}% Rutas por defecto 
    % --- Imagen adaptativa ---
    % Registros de dimensión definidos una sola vez
    \newdimen\imgcap
    \newdimen\imgmin
    \newdimen\imgmax
    \newdimen\imgremain
    \newdimen\imgh
    \newdimen\imgneed
    
    \newcommand{\imagenfit}[3]{%
      \addvspace{10pt}% separación antes
      \par\begingroup
        % Parámetros de composición (asignaciones, no \newdimen)
        \imgcap=1\baselineskip            % alto estimado de título+fuente
        \imgmin=0.15\textheight
        \imgmax=0.15\textheight
        \imgremain=\pagegoal
        \ifdim\pagegoal=\maxdimen \imgremain=0pt\fi
        \advance\imgremain by -\pagetotal
        \advance\imgremain by -\imgcap
        % Decisión de altura destino
        \ifdim\imgremain<\imgmin
          \newpage
          \imgh=0.20\textheight
        \else
          \imgh=\imgremain
          \ifdim\imgh>\imgmax \imgh=\imgmax \fi
        \fi
        % Reservar espacio para evitar cortes del bloque completo
        \imgneed=\imgh \advance\imgneed by \imgcap
        \Needspace{\imgneed}
        % Cuerpo
        \noindent\begin{minipage}{\textwidth}
          \centering
          \captionof{figure}{\textemdash\ #2}\par\vspace{0.3em}% título arriba
          \includegraphics[width=\linewidth,height=\imgh,keepaspectratio]{\detokenize{#1}}\par
          \vspace{0.3em}\small\textit{Fuente: }#3% fuente abajo
        \end{minipage}%
      \endgroup
      \par\addvspace{10pt}%
    }

%------------ Bloque de RECUADRO ESPECIAL -------------%
    \newenvironment{specialblock}[1]{%
      \begin{quote} % crea sangría a izquierda y derecha
      \small % texto más pequeño
      \noindent\textbf{#1}\\[-0.3em] % título en negrita
      \rule{\linewidth}{0.4pt}\\[0.6em]}{
      \end{quote}}
      
%-------------------------- Portada -----------------------%
\newcommand{\oldtitle}[1]{\def\OldTitle{#1}}
\newcommand{\changerationale}[1]{\def\ChangeWhy{#1}}

\newcommand{\changebox}{%
\begin{tcolorbox}[title={Historial de cambios del tema}]
\textbf{Título anterior:} \OldTitle\par
\textbf{Motivo del cambio:} \ChangeWhy
\end{tcolorbox}
}

%-------------------------- Author Font -----------------------%
    \newcommand{\authorfont}[1]{{\fontfamily{EBGaramond-TLF}\selectfont #1}}

%-------------------------- Anexos -----------------------%
    \makeatletter
    \newcounter{anexo}
    \renewcommand{\theanexo}{\Roman{anexo}}
    \newcommand{\anexo}[1]{%
      \clearpage
      \phantomsection
      \refstepcounter{anexo}%
      \noindent\textbf{Anexo \theanexo.- }#1\par\medskip
      \addcontentsline{stc}{section}{Anexo \theanexo.- #1}%
    }
    \makeatother

    %-------------------------- Lista de figuras (personal) -----------------------%
\makeatletter
\newcommand{\MyListOfFigures}{%
  \clearpage
  \phantomsection
  {\centering\bfseries\Large Índice de figuras\par}%
  \medskip
  % Entrada en nuestro índice de contenido (stc)
  \addcontentsline{stc}{section}{Índice de figuras}%
  % Recuperación del listado (sin cabecera propia)
  \begingroup
    \parindent=0pt\parskip=2pt
    \@starttoc{lof}%
  \endgroup
}
\makeatother

%-------------------------- Bibliografía -----------------------%
\makeatletter
% Contador para las referencias
\newcounter{myref}
% Cabecera de la bibliografía, entra en el índice 'stc'
\newcommand{\MyBibliography}{%
  \clearpage
  \phantomsection
  {\centering\bfseries\Large Bibliografía y referencias\par}%
  \medskip
  \addcontentsline{stc}{section}{Bibliografía y referencias}%
  \setcounter{myref}{0}%
}
\newcommand{\MyBibItem}[4]{%
  \refstepcounter{myref}%
  \noindent[\themyref]\ #1.\ \emph{#2}. #3. #4\par
}
\makeatother

%--------- Establecer cuatro niveles de índice --------%
    \setcounter{tocdepth}{4}   
    \setcounter{secnumdepth}{4}% numera hasta \paragraph

%%%%%%%%%%%%%%%%%%%%%%%%%%%%%%%%

\makeatletter

    %--------- Interlineado Global --------%
    \edef\GlobalBaseLineStretch{\baselinestretch}
    
    %--------- Espaciado de seccinones --------%
    \renewcommand\section{\@startsection{section}
      {1}{\z@}
      {2.5ex \@plus 1ex \@minus .2ex} % espacio antes
      {1.5ex \@plus .2ex}             % espacio después
      {\normalfont\Large\bfseries}    % formato del título
    }
    \renewcommand\subsection{\@startsection{subsection}{2}{\z@}
      {3ex \@plus 0.8ex \@minus .2ex}
      {1.5ex \@plus .2ex}
      {\normalfont\large\bfseries}
    }
    \renewcommand\subsubsection{\@startsection{subsubsection}{3}{\z@}
      {3ex \@plus 0.5ex \@minus .2ex}
      {1ex \@plus .2ex}
      {\normalfont\normalsize\bfseries}
    }
    \renewcommand\paragraph{\@startsection{paragraph}{4}{\z@}%
    {-2ex \@plus -0.5ex \@minus -.2ex}% espacio antes
    {0.8ex \@plus .2ex}% espacio después
    {\normalfont\normalsize\itshape}}% ← cursiva en lugar de negrita

    %--------- Ecuaciones --------%   
    % Variables globales para almacenar títulos y notas
    \gdef\leq@current@section{}%
    \gdef\leq@current@subsection{}%
    \gdef\leq@last@printed@section{}%
    \gdef\leq@last@printed@subsection{}%
    
    % Comando para añadir nota (invisible en el cuerpo)
    \newcommand{\eqnote}[1]{%
      \addcontentsline{leq}{leqnote}{#1}%
    }
    
    % Comando para ecuaciones
    \newcommand{\eqblock}[2]{%
      % Verificar si necesitamos imprimir el título de sección
      \ifx\leq@current@section\leq@last@printed@section\else
        \addcontentsline{leq}{leqsection}{\leq@current@section}%
        \global\let\leq@last@printed@section\leq@current@section
        \gdef\leq@last@printed@subsection{}% Reset subsección al cambiar sección
      \fi
      % Verificar si necesitamos imprimir el título de subsección
      \ifx\leq@current@subsection\leq@last@printed@subsection\else
        \ifx\leq@current@subsection\@empty\else
          \addcontentsline{leq}{leqsubsection}{\leq@current@subsection}%
          \global\let\leq@last@printed@subsection\leq@current@subsection
        \fi
      \fi
      % Ecuación
      \refstepcounter{eqnum}%
      \begin{equation*}
        #1\tag{E.\theeqnum}
      \end{equation*}%
      \addcontentsline{leq}{leqt}{[E.\theeqnum] -- #2}%
      \addcontentsline{leq}{leqm}{#1}%
    }
    
    %%% ----- Índice de ecuaciones ----------%%%
    % Tamaño para []
    \newcommand{\leqIndexFont}{\normalsize}
    \newcommand{\leqTitleFont}{\tiny}
    \newcommand{\leqNoteFont}{\tiny}
    % Cabecera de sección 
    \newcommand*\l@leqsection[2]{%
      \addvspace{25pt}% Mayor separación antes de nueva sección
      \noindent\leqIndexFont
      \makebox[\linewidth]{\rule{.38\linewidth}{0.4pt}\ \ \textbf{  \MakeUppercase{#1}}\ \ \rule{.38\linewidth}{0.4pt}}%
      \par\vspace{-10pt}
    }
    % Cabecera de subsección 
    \newcommand*\l@leqsubsection[2]{%
      \par\addvspace{15pt}%
      \noindent\leqIndexFont #1
      \par\vspace{-7pt}
    }
    % Título de ecuación 
    \newcommand*\l@leqt[2]{%
    \par\addvspace{10pt}%
      \par\noindent\leqTitleFont #1\par
    }
    % Miniatura de ecuación
    \newcommand*\l@leqm[2]{%
      \vspace{0.3\baselineskip}%
      \par\scriptsize\noindent\hspace*{1.5em}\(#1\)\par%
    }
    % Nota de ecuación 
    \newcommand*\l@leqnote[2]{%
      \vspace{0.2\baselineskip}%
      \par\noindent\leqNoteFont\hspace*{1.5em}\textbf{Nota.--} #1\par\vspace{0.5\baselineskip}%
    }
    % Generación del índice
    \newcommand{\mylistofequations}{%
      \clearpage
      \phantomsection
      % Título visual de la lista (ajústalo a tu gusto):
      {\centering\bfseries\Large Índice de ecuaciones\par}
      % Entrada en nuestro índice 'stc':
      \addcontentsline{stc}{section}{Índice de ecuaciones}%
      % Llamada a tu lista real (usa el nombre exacto de tu macro):
      \@starttoc{leq}%
    }
    
    %--------- Captura de títulos de sección --------%
    \let\leq@original@section\section
    \let\leq@original@subsection\subsection
    % Flag para saber si estamos en una sección asterisco
    \newif\ifLeqStarSection
    \LeqStarSectionfalse
    
    % Redefinir \section CON CONTROL DE ASTERISCO
    \renewcommand{\section}{%
      \@ifstar{\leq@section@star}{\@dblarg\leq@section@nostar}%
    }
    \def\leq@section@star#1{%
      \global\LeqStarSectiontrue
      \gdef\leq@current@section{#1}%
      \gdef\leq@current@subsection{}%
      \leq@original@section*{#1}%
      \global\LeqStarSectionfalse
    }
    \def\leq@section@nostar[#1]#2{%
      \global\LeqStarSectionfalse
      \gdef\leq@current@section{#2}%
      \gdef\leq@current@subsection{}%
      \leq@original@section[#1]{#2}%
    }
    % Redefinir \subsection
    \renewcommand{\subsection}{%
      \@ifstar{\leq@subsection@star}{\@dblarg\leq@subsection@nostar}%
    }
    \def\leq@subsection@star#1{%
      \global\LeqStarSectiontrue
      \gdef\leq@current@subsection{#1}%
      \leq@original@subsection*{#1}%
      \global\LeqStarSectionfalse
    }
    \def\leq@subsection@nostar[#1]#2{%
      \global\LeqStarSectionfalse
      \gdef\leq@current@subsection{#2}%
      \leq@original@subsection[#1]{#2}%
    }

    % ----- Restauración de valores Globales de estilo ----------%
    % Flags
    \newif\ifInFootnote
    \InFootnotefalse
    \pretocmd{\@footnotetext}{\InFootnotetrue}{}{}
    \apptocmd{\@footnotetext}{\InFootnotefalse}{}{}
    % Profundidad de listas (soporta anidadas)
    \newcount\MyListDepth \MyListDepth=0
    \AtBeginEnvironment{la}{\advance\MyListDepth by 1}
    \AfterEndEnvironment{la}{\advance\MyListDepth by -1}
    \AtBeginEnvironment{lnum}{\advance\MyListDepth by 1}
    \AfterEndEnvironment{lnum}{\advance\MyListDepth by -1}
    \AtBeginEnvironment{itemize}{\advance\MyListDepth by 1}
    \AfterEndEnvironment{itemize}{\advance\MyListDepth by -1}
    \AtBeginEnvironment{enumerate}{\advance\MyListDepth by 1}
    \AfterEndEnvironment{enumerate}{\advance\MyListDepth by -1}
    % Restauración diferida: hazlo al INICIO del próximo párrafo real
    \newif\if@pendingRestore
    \@pendingRestorefalse
    \newcommand{\RequestRestore}{\global\@pendingRestoretrue}
    \everypar{%
      \if@pendingRestore
        \global\@pendingRestorefalse
        \ifInFootnote\else
          \ifnum\MyListDepth=0
            \setlength{\parskip}{\GlobalParskip}%
            \setstretch{\GlobalBaseLineStretch}%
          \fi
        \fi
      \fi
    }
    % Al final de bloques:
    \AfterEndEnvironment{equation}{\RequestRestore}
    \AfterEndEnvironment{equation*}{\RequestRestore}
    \AfterEndEnvironment{la}{\RequestRestore}
    \AfterEndEnvironment{lnum}{\RequestRestore}
    % Al final de macros:
    \apptocmd{\eqblock}{\RequestRestore}{}{}
    \apptocmd{\imagenfit}{\RequestRestore}{}{}
    
\makeatother

% ===================== ToC propio con 4 niveles (único bloque) =====================
\makeatletter

% --- Escritores: añadimos una línea al .stc en cada cabecera numerada ---
% Guardamos las versiones actuales para llamar al comportamiento normal
\let\MyTOC@orig@section\section
\let\MyTOC@orig@subsection\subsection
\let\MyTOC@orig@subsubsection\subsubsection
\let\MyTOC@orig@paragraph\paragraph

% SECTION
\renewcommand{\section}{%
  \@ifstar{\MyTOC@section@star}{\@dblarg\MyTOC@section@nostar}%
}
\def\MyTOC@section@star#1{%
  \MyTOC@orig@section*{#1}% sin entrada al .stc
}
\def\MyTOC@section@nostar[#1]#2{%
  \MyTOC@orig@section[{#1}]{#2}% comportamiento normal (escribe al .toc)
  % Escribimos también nuestra línea al .stc (texto con \numberline)
  \addcontentsline{stc}{section}{\protect\numberline{\thesection}#2}%
}

% SUBSECTION
\renewcommand{\subsection}{%
  \@ifstar{\MyTOC@subsection@star}{\@dblarg\MyTOC@subsection@nostar}%
}
\def\MyTOC@subsection@star#1{%
  \MyTOC@orig@subsection*{#1}%
}
\def\MyTOC@subsection@nostar[#1]#2{%
  \MyTOC@orig@subsection[{#1}]{#2}%
  \addcontentsline{stc}{subsection}{\protect\numberline{\thesubsection}#2}%
}

% SUBSUBSECTION
\renewcommand{\subsubsection}{%
  \@ifstar{\MyTOC@subsubsection@star}{\@dblarg\MyTOC@subsubsection@nostar}%
}
\def\MyTOC@subsubsection@star#1{%
  \MyTOC@orig@subsubsection*{#1}%
}
\def\MyTOC@subsubsection@nostar[#1]#2{%
  \MyTOC@orig@subsubsection[{#1}]{#2}%
  \addcontentsline{stc}{subsubsection}{\protect\numberline{\thesubsubsection}#2}%
}

% PARAGRAPH
\renewcommand{\paragraph}{%
  \@ifstar{\MyTOC@paragraph@star}{\@dblarg\MyTOC@paragraph@nostar}%
}
\def\MyTOC@paragraph@star#1{%
  \MyTOC@orig@paragraph*{#1}%
}
\def\MyTOC@paragraph@nostar[#1]#2{%
  \MyTOC@orig@paragraph[{#1}]{#2}%
  % Si no numeras los \paragraph, cambia \theparagraph por texto plano sin \numberline
  \addcontentsline{stc}{paragraph}{\protect\numberline{\theparagraph}#2}%
}

% --- Impresor del índice propio (lee el .stc) ---
\newcommand{\MyTOC}{%
  \par\bigskip
  {\centering\bfseries\Large Índice de contenido\par}%
  \medskip
  \begingroup
    \parindent=0pt\parskip=2pt
    \@starttoc{stc}%
  \endgroup
}

\makeatother
% ===================== fin ToC propio =====================


%%%%%%%%%%%%%%


% --- Datos del tema ---
\title{Tema 3.A.21 -- La Teoría del Equilibrio General\\
\large }
\author{Marcelo Ortega}
\date{Febrero 2026}
\newcommand{\TemaCode}{3A21}

\oldtitle{Tema 3.A.21. La Teoría del Equilibrio General}
\changerationale{
    Sin cambios en el título. Dada la existencia de un único tema expresamente para el equilibrio general en el temario, se recomienda no acotar la exposición únicamente al modelo de equilibrio general competitivo.
}

\begin{document}


\maketitle
\changebox

\newpage

% --- Página de resumen y modificaciones ---
\puntosclave{ %% Escribir puntos clave Apuntes CECO
    
    }
\vspace{1cm}

\modificaciones{
    \textbf{Consultar:}
    \begin{lnum}
        \item \href{https://competitionandappropriation.econ.ucla.edu/wp-content/uploads/sites/95/2020/12/DebreuMathematicalTheoryofEconomicEquilibrium.pdf?utm_source=chatgpt.com}{\textit{Four Aspects of the Mathematical Theory of Economic Equilibrium}. DEBREU (1975)}
        \item \href{https://scispace.com/pdf/an-introduction-to-general-equilibrium-with-incomplete-asset-192izl9ky2.pdf}{\textit{An introduction to General Equilibirum with incomplete asset's markets.} GEANAKOPLOS (1989)}
        \item \href{https://www.ief.es/docs/destacados/publicaciones/papeles_trabajo/2002_35.pdf}{\textit{Simulación de Políticas Públicas: los modelos de equilibrio general aplicados} GOMEZ-PLANA (2002)}
    \end{lnum}

    \textcolor{red}{OJO ->} Convendría incluir la crítica (o al menos lo que comenta) PICKETTY en el libro de \textit{Igualdad}: ¿Cómo podemos compatibilizar la Teoría del Equilibrio General dentro de un marco en el cual existe un agente económico, cuyo papel en la Economía es sin duda destacable, que fija un conjunto relativamente amplío de los \textit{precios} (en sentido abstracto, pero también técnico) mediante Presupuestos Generales? Resulta, como mínimo, un tanto confuso sostener esta Teoría, o sus posibles implicaciones, observando como funciona en la práctica una \textit{gran porción} de la ésta.

    También, \textbf{ojo ->} Considera la economía desde una óptica de flujo y estática, descarta su concepción dinámica y de stock.

    -> \textit{\textbf{Economía monetaria de producción}}. KEYNES
    }

\newpage
\MyTOC
\newpage

\mylistofequations
\clearpage

\MyListOfFigures
\clearpage


\pagenumbering{arabic}
\setcounter{page}{1}


\section{Introducción}\label{intro}


\subsection{Contextualización}\label{context}


\subsubsection{Relevancia}\label{importance}


\subsubsection{Historia}\label{history}



\subsubsection{Problemática}\label{problem} 




\subsection{Estructura de la exposición}\label{structure}





\clearpage
\section{Equilibrio General: Economía pura de intercambio}
\puntosclave{
    Para WALRAS, Los precios son la variable que se ajusta para llegar al equilibrio.
    
	El equilibrio general se alcanza simultáneamente en todos los mercados y para todos los agentes a través de la interacción de sus decisiones de oferta y demanda.
    
	Si el vector de precios no es el de equilibrio, no se produce ningún intercambio en la economía. El subastador ajusta el vector de precios por medio de un proceso de tanteo hasta lograr el vaciado de todos los mercados.
    
	Puede demostrarse la existencia del equilibrio general. Su estabilidad y unicidad solo pueden demostrarse bajo supuestos restrictivos.
	La asignación de equilibrio alcanzada es eficiente en sentido de Pareto.
}
	
El estudio/análisis de los mercados desde una perspectiva parcial exige que nos imaginemos que el consumidor, dotado de una cantidad concreta de un bien numerario, dispone de forma óptima la cantidad de éste bien que estaba dispuesto a intercambiar por el bien o servicio que se intercambia en dicho merado, dada la estructura de sus preferencias que, por simplicidad y para caracterizarlo, consideramos cuasi-lineales. De esta forma, la curva de demanda del mercado se construye mediante la agregación horizontal de las valoraciones marginales de todos los consumidores presentes y, por otro lado, la curva de oferta del mercado existía o no en función de la estructura de competencia que subyace en el mercado. Resultando todo esto en un equilibrio único, una multiplicidad de equilibrios o la inexistencia de equilibrio. 

El análisis del Equilibrio General es, sin pérdida de generalidad, una alteración tanto metodológica como sustantiva, no muy profunda, de esta situación.

Metodológicamente, el enfoque de equilibrio general tiene dos características centrales:
\begin{lnum}
    \item Primero, concibe la economía como un sistema cerrado e interrelacionado en el que debemos determinar simultáneamente los valores de equilibrio de todas las variables de interés. Así, cuando evaluamos los efectos de una perturbación en el entorno económico, es necesario volver a calcular los niveles de equilibrio del conjunto completo de variables endógenas de la economía. Esto contrasta con el enfoque de equilibrio parcial, en el que el impacto sobre las variables endógenas no directamente relacionadas con el problema en cuestión se desatiende explícita o implícitamente; y,
    \item Una segunda característica central del enfoque de equilibrio general es que aspira a reducir el conjunto de variables tomadas como exógenas a un pequeño número de realidades físicas (por ejemplo, el conjunto de agentes económicos, las tecnologías disponibles, las preferencias y las dotaciones físicas de bienes de los distintos agentes).
\end{lnum}

Desde un punto de vista sustantivo, la teoría del equilibrio general tiene un significado más específico: es una teoría de la determinación de precios y cantidades de equilibrio en un sistema de mercados perfectamente competitivos. 


En sentido estricto, introdujimos un caso particular del modelo de equilibrio general en el Capítulo 10. Allí llevamos a cabo un análisis de equilibrio y de bienestar de mercados perfectamente competitivos bajo el supuesto de que los consumidores tenían preferencias cuasilineales. En ese contexto, las funciones de demanda de los consumidores no presentan efectos riqueza (salvo para un único bien, denominado numerario); como consecuencia, el análisis de un solo mercado (o de un pequeño grupo de mercados) podía desarrollarse de un modo comprensible como el análisis tradicional de equilibrio parcial. Una buena parte de lo que hacemos en la Parte IV puede verse como un intento de extender las ideas del Capítulo 10 a un mundo en el que los efectos riqueza son significativos. La motivación principal de ello es el aumento de realismo que conlleva. Para hacer un uso práctico del análisis de equilibrio en el estudio del funcionamiento de una economía en su conjunto, o para evaluar intervenciones de política que afectan simultáneamente a un gran número de mercados, los efectos riqueza, una fuente primaria de interconexiones entre mercados, no pueden ser ignorados, y por lo tanto el enfoque de equilibrio general es esencial.

Aunque el conocimiento del material tratado en el Capítulo 10 no es un requisito estricto para la Parte IV, recomendamos no obstante encarecidamente que se estudie, especialmente las Secciones 10.B a 10.D. Constituye una introducción a las cuestiones principales y proporciona un ejemplo sencillo y analíticamente muy útil. Veremos en los distintos capítulos de la Parte IV que un buen número de los resultados importantes establecidos en el Capítulo 10 para el caso cuasilineal se extienden al caso de preferencias generales. Pero muchos otros no. Para entender por qué puede ser así, recuérdese de los Capítulos 4 y 10 que un grupo de consumidores con preferencias cuasilineales (con respecto al mismo numerario) admite la existencia de un consumidor representativo (normativo). Esta es una restricción poderosa sobre el comportamiento de la demanda agregada que no estará disponible en los entornos más generales que estudiamos aquí.

Es importante señalar que, en relación con el análisis llevado a cabo en la Parte III, incurrimos en un coste para lograr la tarea que el equilibrio general se propone realizar: los supuestos de comportamiento tomador de precios y de cotización universal de precios —esto es, la existencia de mercados para toda mercancía relevante (con la implicación de información simétrica)— están presentes en casi toda la teoría estudiada en la Parte IV. Así, en muchos aspectos, no profundizamos tanto como lo hicimos en la Parte III en el microanálisis de los mercados, de los fallos de mercado y de la interdependencia estratégica de los agentes del mercado. El compromiso en la estructura conceptual entre las Partes III y IV refleja, en cierto sentido, el estado actual de la frontera de la investigación microeconómica.

El contenido de la Parte IV se organiza en seis capítulos.

El Capítulo 15 presenta una discusión preliminar. Su propósito principal es ilustrar las cuestiones que conciernen a la teoría del equilibrio general mediante tres ejemplos sencillos: la economía de caja de Edgeworth con dos consumidores; la economía con un consumidor y una empresa; y el modelo de economía pequeña y abierta.

Los Capítulos 16 y 17 constituyen el núcleo del análisis formal de la Parte IV. El Capítulo 16 presenta la estructura formal del modelo de equilibrio general e introduce dos conceptos centrales de la teoría: las nociones de optimalidad de Pareto y de equilibrio con toma de precios (y, en particular, equilibrio walrasiano). El capítulo está dedicado al examen de la relación entre estos dos conceptos. El énfasis es, por tanto, normativo, centrado en las propiedades de bienestar de los equilibrios con toma de precios. El núcleo del capítulo se ocupa de la formulación y demostración de los dos teoremas fundamentales de la economía del bienestar.

En el Capítulo 17, el énfasis se pone, en cambio, en las propiedades positivas (o descriptivas) de los equilibrios walrasianos. Estudiamos una serie de cuestiones relativas al poder predictivo de la teoría walrasiana, incluida la existencia, la unicidad local y global, y el comportamiento de estática comparativa de los equilibrios walrasianos.

Los Capítulos 18 a 20 exploran extensiones del análisis básico presentado en los Capítulos 16 y 17. El Capítulo 18 abarca una serie de temas cuyos orígenes se encuentran en la teoría normativa o en la teoría cooperativa de juegos; estos temas comparten la característica de que proporcionan una mirada más profunda a los fundamentos de los equilibrios con toma de precios mediante la explotación de propiedades derivadas de la naturaleza masiva de los mercados. Estudiamos el importante teorema de equivalencia del núcleo, examinamos con mayor detalle la idea de los equilibrios walrasianos como el límite de los equilibrios no cooperativos a medida que los mercados se hacen grandes (un tema ya abordado en la Sección 12.F) y presentamos dos caracterizaciones normativas de los equilibrios walrasianos: una en términos de ausencia de envidia (o anonimato) y otra en términos de un principio de productividad marginal. El Apéndice A del Capítulo 18 ofrece una breve introducción a la teoría cooperativa de juegos.

El Capítulo 19 aborda la modelización de la incertidumbre en un contexto de equilibrio general. La capacidad de hacerlo de un modo teóricamente satisfactorio ha sido una de las historias de éxito de la teoría del equilibrio general. Aquí se introducen y estudian los conceptos de bienes contingentes, equilibrio de Arrow-Debreu, comercio secuencial (en un entorno de dos períodos), equilibrio de Radner, arbitraje, equilibrio con expectativas racionales y mercados incompletos. El capítulo proporciona un vínculo natural con la teoría moderna de las finanzas.

El Capítulo 20 considera la aplicación de la teoría general a economías competitivas dinámicas (pero sin incertidumbre) y estudia asimismo una serie de cuestiones específicas de este entorno. Se introducen nociones como la impaciencia, la eficiencia dinámica y la maximización miope frente a la maximización de la utilidad global. El capítulo analiza en primer lugar economías dinámicas con consumidor representativo (incluido el modelo de Ramsey-Solow), luego generaliza al caso de un número finito de consumidores y concluye con una breve presentación del modelo de generaciones solapadas. En el proceso, exploramos una amplia gama de comportamientos dinámicos. El capítulo proporciona un vínculo natural con la teoría macroeconómica.

Los clásicos modernos de la teoría del equilibrio general son Debreu (1959) y Arrow y Hahn (1971). Estos textos ofrecen una discusión adicional de los temas tratados aquí. Para extensiones, recomendamos la cobertura enciclopédica de Arrow e Intriligator (1981, 1982, 1986) y Hildenbrand y Sonnenschein (1991). Véase también el relato más reciente en forma de manual de Ellickson (1993). El análisis de equilibrio general tiene una dimensión aplicada muy importante que no abordamos en este libro, pero que explica en buena medida la importancia de la teoría. Para una revisión, recomendamos Shoven y Whalley (1992).

\subsection{Desarrollo formal}
Como es evidente, la teoría walrasiana de los mercados es muy ambiciosa. Intenta nada menos que predecir el vector completo de consumos y producciones para un economía: (i) la lista de bienes, $L\gg 0$; (ii) el número de consumidores, $I\gg 0$; y, (iii) el número de productores, $J\gg 0$. 

\textbf{NOTA.-} Es necesario hacer un briefing de todo el desarrollo teórico que precede al modelo que expondremos. En particular, cabe hacer mención a que presentaremos un caso particular del Equilibrio General Competitivo en una economía de producción completamente descentralizada. Sin embargo, un modelo primitivo sería el Equilibrio General Competitivo en una economía de intercambio puro, donde realmente se definen de forma conceptual las funciones que dan lugar al Equilibrio. Además, aunque se mencionará cuando se impongan, los supuestos que se presentan correposponden también a un caso particular, con el fin de hacerlo simple y sencillo sin pérdida de generalidad. 


¿Por qué no importa esta simplificación?

- Simplicidad analítica
- Siempre es posible reducir una economía con producción a una economía de intercambio en la que, en efecto, los consumidores intercambian insumos factoriales y luego realizan producción doméstica utilizando una tecnología de rendimientos constantes libremente disponible. El exceso de demanda agregado en esta economía de intercambio derivada para los insumos factoriales combina elementos tanto del consumo como de la producción y bien podría satisfacer la propiedad de sustitución bruta.⁴⁷
- Todos los resultados obtenidos pueden extenderse a una economía descentralizada con producción (Modelo de ARROW-DEBREU). Cuando esto no sea posible, así se explicitará.


\subsubsection{Supuestos}
Para ello, parte de los siguientes supuestos.

\paragraph{Fundamentos de la Economía}
Relativos a los fundamentales de la economía, se se asume que:
\begin{lnum}
    \item \textbf{Consumidores}. Cada consumidor viene caracterizado por un conjunto de consumo $X_i\subset\mathbb R^L$ y una relación de preferencias $\succsim_i$ definida dentro de este conjunto de consumo. Asumimos que estas preferencia son racionales continuas, estríctamente convexas, y localmente insaciables;
    \item \textbf{Dotaciones iniciales}. Asumimos también que cada consumidor dispone de una dotación inicial de biene $\omega_i\in\mathbb R^L$; y, 
    \item \textbf{Productores}. Al tratarse de una economía pura de intercambio, la única actividad de producción posible es la eliminación gratuita. Formalmente, $J=1$ y $Y_1=-\mathbb R^L$.
\end{lnum}

\paragraph{Institucionales}
Por lo que respecta a los supuestos institucionales, supondremos que en el equilibrio se fija un precio para cada bien, incluidos aquellos que no se comerciarán, y que no existen restricciones informativas. 

\paragraph{Conductuales}
\textcolor{red}{Por último, desde una perspectiva conductual de los agentes, suponemos que todos los agentes son precio-aceptantes. Esto requiere la consideración de una economía \textbf{grande} o \textbf{atomizada} y sus implicaciones son cruciales para el equilibrio: sin esta propiedad, podrían existir no-convexidades en los conjuntos de eleción, lo que podría implicar la inexistencia de un equilibrio. 

Puede entonces resultar contraintuitivo que a lo largo de la exposición supongamos un conjunto dado de agentes $I$ (consumidores pero también extensible a empresas). La razón es que se puede justificar que, con economías suficientemente grandes, incluso en presencia de no-convexidades en las preferencias las conclusiones que nos disponemos a presentar se mantienen}.\footnote{%
    Para una demostración sencilla e intuitiva en el caso de una economía de intercambio pura: 

Supongamos que consideramos economías cuyos consumidores tienen características (preferencias y dotaciones) que se agrupan en $I$ tipos dados, con $r$ consumidores de cada tipo (es posible una generalización a números desiguales por tipo, $r_i$). 

El conjunto de consumidores está formado por $r$ réplicas de un conjunto básico de referencia de consumidores. Además, una asignación denotada por $(x_1,\ldots,x_I)$ se entiende ahora como que especifica que cada consumidor del tipo $i$ consume $x_i$ (de modo que el conjunto de consumidores del tipo $i$ consume $r\cdot x_i$). Observamos entonces que el análisis y los resultados presentados hasta este punto no se modifican por esta reinterpretación; simplemente no dependen en modo alguno del parámetro $r$. De esta manera, podemos concluir informalmente que la teoría cubre casos con un número arbitrariamente grande, incluso infinito, de consumidores; en particular, vemos que cualquier equilibrio del modelo que presentaremos es un equilibrio de la economía $r$-réplica (para cualquier entero $r\ge 1$).

No obstante, esto no resuelve el problema: La capacidad de interpretar el modelo y los resultados de una manera que sea completamente independiente del número de consumidores depende de manera crucial del supuesto de convexidad de las preferencias. Sin este supuesto, no está justificado descuidar asignaciones que asignen cestas de consumo diferentes a distintos consumidores del mismo tipo. No obstante, esta reinterpretación sugiere una observación interesante: la replicación puede de hecho ayudar en el análisis de economías con no convexidades, en el sentido de que un aumento en el tamaño de la economía (en términos del número de réplicas) puede ayudar a asegurar la existencia de un equilibrio. En efecto, si la economía es lo suficientemente grande, entonces la existencia de un equilibrio está asegurada, o casi, incluso si las preferencias no son convexas.
}

\subsubsection{Caracterización del Equilibrio}
Bajo este marco formal, podemos entonces definir el Equilibrio General como la conjunción de las siguientes condiciones:

\eqblock{
\begin{aligned}
    &\text{Se una economía definida por:}\qquad \left(\{(X_i,\succsim_i)\}_{i=1}^I,\;\{\omega_i\}_{i=1}^I,\;Y_1\right)\\[2em]
    &((\mathbf x^*,\mathbf y^*),\;\mathbf p\in\mathbb R^L)\equiv \underset{\text{\scriptsize Walrasiano}}{\text{Equilibrio}}\iff\;
    \left\{\begin{aligned}
        &\text{(1)}&& y_1^*\leq 0,\;\mathbf p\cdot y_1^*=0\quad \text{y,}\;\;\mathbf p\geq0\\[1em]
        &\text{(2)}&& \forall i,\quad\;x_i^*=x_i(\mathbf p,\;\mathbf p\cdot \omega_i)\\[1em]
        &\text{(3)}&& \sum_{i=1}^Ix_i^*-\sum_{i=1}^I\omega_i=y_1^*
    \end{aligned}\right.
\end{aligned}
}{Definición de EGC: Condiciones de Equilibrio. Equilibrio General: Ausencia de fallos de mercado}

La segunda y tercera condición son inmediatas, mientras que la primera implica que en equilibrio las dotaciones eliminadas deben ser negativas o nulas (restricción tecnológica) y necesariamente ausentes de valor. Es decir, o bien no se elimina ninguna dotación, o las dotaciones que se eliminan no tienen precio (son gratuitas) y, en consecuencia, no reflejan la valoración del consumidor (esto es imprescindible puesto que si tuviesen valoración, se reflejaría en el precio y consecuentamente no se eliminarían).

Además, en una economía pura de intercambio que cumple los supuestos sobre preferencias, podemos ir un paso más allá saber que cualquier vector de precios no-negativo es un vector de precios de Equilibrio General si y solo si cumple:

\eqblock{
\begin{aligned}
    \boxed{\mathbf p\equiv\text{EGC} \iff \sum_i(x_i(\mathbf p,\;\mathbf p\cdot \omega_i)-\omega_i)\leq0}
\end{aligned}
}{}

Esta condición se desprende automáticamente de la primera y tercera conclusión. No obstante, a efectos del análisis formal, resulta sumamente útil poder expresar los equilibrios como las soluciones de un sistema de ecuaciones. Objetivo que perseguiremos en adelante, aspirando a ser muy concretos pero debiendo recurrir a la imposición de supuestos más fuertes para simplificar el análisis.

\paragraph{Función de Exceso de Demanda Agregada}
Para ello introducimos la noción de \textbf{función de exceso de demanda}. Definimos la Función de Exceso de Demanda del consumidor $i$ dado el vector de precios $\mathbf p$, como:

\eqblock{
\begin{aligned}
    &\text{FED Individual: }\;\underset{\text{\scriptsize Se puede interpretar como: \textit{Lo que desearía menos lo que inicialmente poseo}}}{\boxed{z_i(\mathbf p)\in\mathbb R^L\equiv \underset{\text{\scriptsize Función de Demanda Walrasiana}}{x_i\left(\mathbf p,\;\mathbf p\cdot\omega_i\right)}-\omega_i}}\\[1em]
    &\text{FED Agregada: }\;\underset{\substack{\\[0.5em]\\\text{\scriptsize donde 
    $\left\{\begin{aligned}
        &z(\mathbf p)>0\sim \text{Exceso de demanda generalizado}\\
        &z(\mathbf p)\leq0 \sim \text{Plétora generalizada de recursos}
    \end{aligned}\right.$}}}{\boxed{z(\mathbf p)\in\mathbb R^L\equiv \sum_iz_i(\mathbf p)}}\\[2em]
    &\text{Utilizando la condición inicial y las nuevas definiciones, tenemos que:}\\[2em]
    &\hspace{4em}\boxed{\mathbf p\in \mathbb R^L\equiv \text{Equilibrio}\;\;\iff\;\;z(\mathbf p)\leq 0}
\end{aligned}
}{Definición}

Por tanto, la demanda neta se define como lo que el consumidor desea consumir menos lo que trae como dotación. De manera general, la Función de Exceso de Demanda establece las condiciones del conjunto de bienes resultantes de la interacción, teniendo cuenta las dotaciones iniciales, las demandas individuales, y la tecnología de producción. 

Nótese que una economía de intercambio puro con preferencias localmente insaciables, si el precio de un bien es positivo ($p_\ell>0$), no puede sobrar en equilibrio, porque habría posibilidad de aumentar beneficios/valor o habría desequilibrio: $p_\ell\geq0,\;\;z_\ell(\mathbf p)\leq 0,\;\; p_\ell z_\ell(\mathbf p)=0$. Entonces, necesariamente cualquier bien $\ell$ puede presentar oferta excesiva ($z_\ell(\mathbf p)<0$) si y solo si este bien es gratuito, $p_\ell=0$.\footnote{%
    Como ejemplo, el bien $\ell$ puede ser un \textit{mal}. Entonces, podríamos esperar que su precio fuese cero, la demanda de los consumidores fuese cero y el exceso de oferta, dado por $z_\ell=\omega_\ell>0$, fuese eliminado mediante la condición de eliminación gratuita.
} Si este bien tiene un precio positivo necesariamente debe vaciarse el mercado pues el valor para algún consumidor sería positivo. 

De hecho, si asumimos preferencias estríctamene monótonas, cualquier equilibrio general debe implicar un vector de precios estríctamente positivo, $\mathbf p\gg0$, ya que, de lo contrario los consumidores demandarían una cantidad infinita de cualquier bien gratuito. Por tanto, obtendríamos:

\eqblock{
\begin{aligned}
    &\text{Si $\succsim_i$ son racionales, transitivas y estríctamente monótonas:}\\[1em]
    &\Longrightarrow\qquad \boxed{\mathbf p=(p_1,\dots p_L)\equiv \text{EGC}\iff z_\ell(\mathbf p)=0,\;\;\forall\ell=1,\dots, L:\quad z(\mathbf p)=0}\\[1em]
    &\text{Y, $z(\mathbf p)$ definida para todo vector $\mathbf p\gg 0$ cumple:}\\[0.5em]
    &\left\{\begin{aligned}
        &\text{(1)}&&z(\cdot )\in\text{Continua}\\
        &\text{(2)}&&z(\cdot )\in\mathcal H(0)\\
        &\text{(1)}&&\mathbf p\cdot z(\mathbf p)=0,\;\;\forall \mathbf p\\
        &\text{(1)}&&\exists s>0:z_\ell(\mathbf p)>-s,\;\;\forall \ell\;\&\;\forall\mathbf p\\
        &\text{(1)}&&\text{Si, }\;\;\mathbf p^n\to \mathbf p,\;\;\text{donde}\;\;\mathbf p\neq0\;\;\&\;\;p_\ell=0\;\;\text{para algún $\ell$},\\
        &&&\hspace{4em}\max\{z_1(\mathbf p^n),\dots ,z_L(\mathbf p^n)\}\to \infty
    \end{aligned}\right.
\end{aligned}
}{}

Con la excepción de la propiedad (v), el resto de propiedades son consecuencia directa de la definición y las propiedades paralelas de las funciones de demanda. La cota o límite $s$ se deriva de la no-negatividad de la demanda, que implica que la oferta neta de un consumidor al mercado de cualquier bien $\ell$ no puede ser mayor que su dotación inicial. La intuición detrás de la quinta propiedad es precisamente que cuando algún úprecio tiende a cero, un consumidor cuyo bienestar/utilidad tiende estríctamente a un límite positivo y cuyas preferencias son estríctamente monótonas, demandará una cantidad cada vez mayor de algún/os bien/es/commodity cuyo/s precio/s tiendan a cero, pero no de todos debido precisamente a la importancia de los precios relativos. 

\paragraph{Ley de WALRAS}
De hecho, podems observar que, debido a las propiedades de la FDE agregada definida sobre el conjunto posible de vectores de precios estríctamente positivos, para verificar si un vector de precios estríctamente superior a cero es un vector de equilibrio (EGC), será suficiente con ver si es un vector de vaciamiento de mercados para todos salvo para uno:

\eqblock{
\begin{aligned}
    &\text{Si $\mathbf p\gg 0$ es un vector de EGC, podemos comprobarlo:}\\[0.5em]
    &\Longrightarrow \mathbf p^*\gg 0:z_1(\mathbf p^*)=\cdots=z_{L-1}(\mathbf p^*)=0\\[0.5em]
    &\underset{\text{\scriptsize como,}}{\Longrightarrow}\;\;\mathbf p^*\cdot z(\mathbf p^*)=\sum_\ell p_\ell^*\cdot z_\ell(\mathbf p^*)=0\;\;\text{y,}\;\; p_L>0\\[0.5em]
    &\Longrightarrow z_L(\mathbf p^*)=0\underset{\substack{\text{\scriptsize denotamos vector FDE}\\\text{\scriptsize para todos los bienes desde 1}\\\text{\scriptsize hasta $L-1$, como}}}{\Longrightarrow}\tilde z(\mathbf p^*)=(z_1(\mathbf p^*),\dots,z_{L-1}(\mathbf p^*))\\[0.5em]
    &\Longrightarrow \underset{\text{\scriptsize Ley de WALRAS: Vaciamiento de mercados}}{\boxed{\mathbf{p}^*\gg 0\equiv \text{EGC}\iff \tilde z(\mathbf p^*)=0}}
\end{aligned}
}{Ley de WALRAS: Vaciamiento de todos los mercados menos uno}

Es decir, como la suma de los excesos de demanda debe ser cero, los mercados necesariamente son interdependientes. Por tanto, si $L-1$ mercados están en equilibrio, el último también lo estará. De manera análoga, si $L-1$ mercados presentan exceso de oferta, el último mercado necesariamente debe presentar exceso de demanda. 

De manera más general, este resultado nos permite obtener el precio de un bien en términos de los demás bienes. Es decir, permite determinar los precios de los bienes en términos relativos sin que se conozcan los valores absolutos de los precios. En conclusión, el equilibrio general competitivo es un conjunto de precios no negativos, asociado a un conjunto de demandas y ofertas de economías domésticas y empresas, tales que cada una de las demandas y ofertas es óptima y los excesos de demanda existentes son nulos (es decir, los planes de los agentes son compatibles y pueden llevarse a cabo).

Ahora ya hemos definido el Equilibrio General Walrasiano y hemos presentado una herramienta analítica que, \textit{a priori}, parece permitirnos simplificar la complejidad que derivaría de un sistema de ecuaciones basado en las condiciones de equilibrio a uno mucho más sencillo que determine el equilibrio de manera global. Pero, ¿qué ventaja nos propociona esto? Evidentemente resulta atractivo, ¿pero que de cierto hay en realidad?

Para contestar a estas cuestiones y ver las implicaciones de Política Económica que nos permite estudiar el modelo de Equilibrio General debemos dividir el análisis en dos enfoques: uno positivo y uno normativo.

\subsection{Implicaciones positivas}
Al estudiar una teoría positiva, la primera pregunta que hay que plantear es: ¿bajo qué condiciones posee el modelo formal una solución? Es decir, ¿es capaz de predecir un resultado definido? Esto se conoce como el problema de la existencia. Conceptualmente, la garantía de la existencia de un equilibrio significa que nuestra noción de equilibrio supera la prueba lógica de la consistencia. Nos dice que el modelo matemático es adecuado para los propósitos para los que ha sido diseñado. Aunque un teorema de existencia difícilmente puede ser el final de la historia, es, en cierto sentido, la puerta que se abre a la casa del análisis.

\subsubsection{Existencia del EGC}
La existencia de un equilibrio walrasiano puede establecerse con considerable generalidad. Para mantener el flujo natural de la exposición, ofrecemos un examen detallado del problema de la existencia para el caso articular de una economía de intercambio puro modeladas por medio de funciones de exceso de demanda. Hemos visto que la función de exceso de demanda $z(\cdot)$ de una economía de intercambio con $\sum_i\omega_i \gg 0$ y preferencias continuas, estrictamente convexas y fuertemente monótonas satisface una serie de propiedades. A continuación demostramos que cualquier función $z(\cdot)$ que satisfaga estas cinco condiciones admite una solución, es decir, un vector de precios $\mathbf p$ tal que $z(\mathbf p)=0$. Al hacerlo, establecemos que existe un equilibrio walrasiano bajo dichas condiciones.

\eqblock{
\begin{aligned}
    &\text{\textbf{Proposición.-}}\qquad \text{Si,}\;\;z(\cdot)
    \left\{\begin{aligned}
        &\text{(1)}&&\mathbb{Dom}\{\mathbf p:\mathbf p\in\mathbb R^L_+\}\\
        &\text{(2)}&&\text{Continua}\\
        &\text{(2)}&&\mathcal H(0)
    \end{aligned}\right\}\Longrightarrow\exists\mathbf p:z(\mathbf p)\leq0\\[2em]
    &\text{\textbf{Prueba:}}\qquad \text{Como, }\;\;z(\cdot)\in\mathcal H(0)\underset{\substack{\text{\scriptsize restringimos nuestra}\\\text{\scriptsize búsqueda al simplex unitario}}}{\Longrightarrow} \Delta=\left\{\mathbf p\in\mathbb R^L_+:\sum_\ell p_\ell=1\right\}\\[0.5em]
    &\over\hspace{4em}
    \begin{aligned}
        \\
        &\left\{\begin{aligned}
            &\text{(1) Definimos:}&&z^+(\cdot)\in\mathbb{Dom}\{\Delta\}: z^+_\ell (\mathbf p)=\max\{z_\ell(\mathbf p), \;0\}&& (\substack{\text{\scriptsize Nótese que $z^+(\cdot)$ es continua y que si}\\\text{\scriptsize $z^+(\mathbf p)\cdot z(\mathbf p)=0$ implica que $z(\mathbf p)\leq0$}})\\[0.5em]
            &\text{(2) Denotamos:}&&\alpha(\mathbf p)=\sum_\ell\left[p_\ell+z_\ell^+(\mathbf p)\right]\Longrightarrow\alpha(\mathbf p)\geq 1,\;\;\forall \mathbf p\\[0.5em]
            &\text{(3) Definimos:}&&f(\cdot):f(\mathbf p)=\left[\frac{1}{\alpha(\mathbf p)}\right](\mathbf p+z^+(\mathbf p))&& \left(\substack{\text{\scriptsize Nótese que intuitivamente se ve que}\\\text{\scriptsize esta función de punto fijo tiende a incrementar}\\\text{\scriptsize el precio de bienes con exceso de demanda}}\right)\\[0.5em]
    \end{aligned}\right.\\[2em]
    &\begin{aligned}
        &\underset{\substack{\text{\scriptsize Por el Teorema}\\\text{\scriptsize del Punto Fijo de}\\\text{\scriptsize BROUWER}}}{\Longrightarrow}&&\exists\mathbf p^*\in\Delta :\mathbf p^*=f(\mathbf p^*)\\
        &\underset{\substack{\text{\scriptsize Por la Ley}\\\text{\scriptsize de WALRAS}}}{\Longrightarrow}&& 0=\mathbf p^*\cdot z(\mathbf p^*)=f(\mathbf p^*)\cdot z(\mathbf p^*)=\left[\frac{1}{\alpha(\mathbf p^*)}\right]z^+(\mathbf p^*)\cdot z(\mathbf p^*)\;\;\Longrightarrow\;\;z^+(\mathbf p^*)\cdot z(\mathbf p^*)=0\\[0.5em] 
        &\Longrightarrow&& z(\mathbf p^*)\leq 0\qquad \qquad\blacksquare
    \end{aligned}
    \end{aligned}
\end{aligned}
}{Teorema de BROUWER: Existencia del Equilibrio General Competitivo. Análisis positivo}

Este desarrollo se puede representar, alternativamente, mediante el gráfico siguiente, restringiendo la economía de intercambio pura a dos bienes:

\imagenfit{BROUWER_Existencia.png}{La proposición del Teorema del Punto Fijo: Existencia del EGC}{\textit{Microeconomic Theory}. MAS-COLELL et al., 1995}

La intuición detrás de la demostración se basa en postular un mecanismo de ajuste que asigna a cada vector de precios otro vector de precios, en función del exceso de demanda al que dé lugar.
a)	Cuando existe exceso de demanda en un mercado, se le asigna un nuevo precio más elevado.
b)	Cuando existe exceso de oferta, se le asigna un precio más bajo.
Este mecanismo es así una aplicación de un conjunto sobre sí mismo: cuando se alcance un vector de precios que anule todos los excesos de demanda, el mecanismo asignará el vector de precios a sí mismo, porque no será necesario variar ningún precio relativo para eliminar ningún exceso de demanda y estaremos ante un punto fijo del sistema que es un equilibrio competitivo. Así, si es posible demostrar la existencia de un punto fijo, se estará también demostrando la existencia de un equilibrio competitivo. De este modo, cuando las funciones de exceso de demanda son i) continuas, ii) homogéneas en grado cero y iii) cumplen la Ley de Walras, se cumple que efectivamente existe un vector de precios que elimina todos los excesos de demanda y, por ende, existe el equilibrio competitivo.

\subsubsection{Unicidad del EGC}
Habiendo establecido las condiciones bajo las cuales se garantiza la existencia de un equilibrio walrasiano, necesitamos saber lístas características de este. Desde una perspectiva teórica, el mejor de todos los mundos posibles es aquel en el que la situación social que se analiza puede formalizarse de una manera que, por un lado, sea muy parsimoniosa (es decir, utilice como insumos únicamente los rasgos más indiscutibles y sólidos de la realidad que se modela) y, por otro, logre predecir un \textbf{resultado único}.\footnote{%
    El modelo walrasiano de competencia perfecta es, en efecto, muy parsimonioso. Esencialmente, intenta dar una explicación teórica completa de una economía utilizando como fundamentos únicamente la lista de bienes, el estado de la tecnología y las preferencias y dotaciones de los consumidores. Sin embargo, el otro lado de la moneda es que la teoría no es completamente determinista: la unicidad de los equilibrios está asegurada solo bajo condiciones especiales.

Si la unicidad global no es alcanzable, la siguiente mejor propiedad es la unicidad local. Decimos que un vector de precios de equilibrio es localmente único, o localmente aislado, si no podemos encontrar otro vector de precios (normalizado) arbitrariamente cercano a él. Si todo equilibrio de una economía es localmente único, decimos que la propiedad de unicidad local se cumple para la economía. La propiedad de unicidad local es de interés porque, si prevalece, entonces puede no ser difícil completar la teoría en cualquier aplicación particular. Por ejemplo, la historia puede haber determinado la región en la que se encuentra el equilibrio (podría ser la región en la que el equilibrio se encontraba antes de un pequeño choque no anticipado a la economía), y en esa región podemos tener un equilibrio único. En este caso, la teoría conserva su poder predictivo, aunque solo localmente. Decimos que una teoría que garantiza la unicidad local de los equilibrios es localmente (a diferencia de globalmente) determinante.
} 

Restringiéndonos al análisis de economías de intercambio puro y manteniendo los supuestos previos, supongamos además que $z(\cdot)$ es continuamente diferenciable.

\paragraph{Unicidad local: Caso general}
La primera pregunta que debemos hacernos es: ¿la teoría walrasiana localmente determinante? Aunque parezca contraintuitivo, la unicidad local es una condición previa y necesaria para la unicidad global, pues un sencillo razonamiento nos lleva a ver que si el equilibrio local es global, el equilibrio alcanzado es el único equilibrio general. La Teoría Microeconómica ha demostrado que:\footnote{%
    \textbf{Definición.-} Una economía regular es aquella en la que los equilibrios tienen matriz Jacobiana/diferencial no singular (que no existen casos degenerados). Por ejemplo, puntos silla o tangencias en equilirbio. En definitivo, la función de exceso de demanda presente en el equilibrio presenta pendiente negativa o positiva estricta.

Aunque no se menciona, la demostrón formal aciseó debe a DEBREU (1970). La medida de Lebesque es nula, por tanto la probabilidad de encontrarse frente a una economía no regular es cero (ojo, siempre y cuando cumpla los supuestos de partida de la demostración).
}

\eqblock{
\begin{aligned}
    &\text{Si,}\qquad \{\tilde z(\mathbf p)\in\mathcal C^n,\quad n\ge2\quad \text{y,}\quad\text{Economía \textbf{regular}}\}\\[1em]
    &\Longrightarrow\qquad  \forall\mathbf p^*:\tilde z(\mathbf p^*)=0,\qquad \exists \varepsilon>0: \;\;(\mathbf p'\neq \mathbf p^*;\;\;\mathbf p'_L=\mathbf p^*_L;\;\;\|\mathbf p'-\mathbf p^*\|<\varepsilon)\Longrightarrow z(\mathbf p')\neq0
\end{aligned}
}

Por tanto, probamos que, bajo una serie de supuestos generales (de los cuales el más sustancial se refiere a la convexidad), existe un equilibrio localmente único. Gráficamente:

\imagenfit{Unicidad_LOCAL.png}{FED para economías de equilibrios múltiples: (izq.) localmente únicos; (dcha.) continuos}{\textit{Microeconomis Theory}. MAS-COLELL et. al., 1995}

Una vez que sabemos que todo equilibrio está localmente aislado, la finitud del número de equilibrios es una consecuencia de la condición de frontera de la FDE (quinta propiedad).\footnote{

Las condiciones de frontera sobre la función de exceso de demanda nos permite demostrar dos elementos cruciales:
\begin{lnum}
    \item El equilibrio no es compatible con precios relativos que sean arbitrariamente cercanos a cero. Es decir, existe un $r > 0$ tal que, si $z(p)=0$ y $p_L = 1$, entonces $1/r < p_t < r$ para todo $t$. La continuidad de $z(\cdot)$ añade a esto el hecho de que el conjunto de vectores de precios de equilibrio es un subconjunto cerrado de $\mathbb{R}^{L-1}$. Pero un conjunto que es cerrado y acotado (es decir, compacto) en $\mathbb{R}^{L-1}$ y discreto (es decir, con todos sus puntos localmente aislados) debe ser necesariamente finito; y,
    \item \textbf{Índice}. La existencia del índice de los vectores de precios de equilibrio ($+1$). Se trata de una propiedad que no analizamos [Véase el Manual] pero que tiene implicaciones interesantes. En primer lugar, que el número de equilibrios de una economía regular es impar. En particular, este número no puede ser cero; por lo tanto, la existencia de al menos un equilibrio es un caso particular de la proposición. Segundo, el concepto de índice proporciona una clasificación de los equilibrios en dos tipos. De hecho, típicamente ocurre que cualquier búsqueda de equilibrios bien comportados (lo que esto significa depende de la aplicación particular) puede restringirse a los equilibrios de índice positivo. Tercero, el resultado del índice tiene implicaciones para la unicidad y la multiplicidad de los equilibrios. Cuarto, parte de la importancia del teorema del índice es que esto es todo lo que podemos esperar derivar sin imponer supuestos adicionales (fuertes).
\end{lnum}

La determinación local genérica de la teoría se extiende a casos con externalidades, impuestos u otras imperfecciones que conducen al fracaso del primer teorema del bienestar. Por otro lado, la finitud del número de equilibrios es una conclusión burda. No es lo mismo que finito signifique tres o que signifique unos pocos millones. Lamentablemente, salvo ir hasta las condiciones de unicidad, no disponemos de ninguna técnica que nos permita refinar nuestras conclusiones. Queremos dejar constancia, sin embargo, de que no debe presumirse que, en toda generalidad, finito signifique pequeño. 
} Hasta este punto, nos hemos concentrado en la determinación de las propiedades generales del modelo de equilibrio walrasiano. Ahora adoptamos un enfoque diferente. Estas son propiedades importantes, pero nos gustaría saber si podemos decir algo más, especialmente con fines predictivos o de estática comparativa. Aunque a estas alturas bien podemos sospechar que la respuesta es negativa sin que sean necesario introducir supuestos especiales para derivar implicaciones más fuertes.

\paragraph{Unicidad global: Supuestos reforzados}
Nos centramos ahora en una propiedad particular muy importante: la unicidad del equilibrio.

Diremos que, en una economía de intercambio, el equilibrio es único si y solo si se cumple la propiedad de \textbf{sustitución bruta} para el exceso de demanda agregado.\footnote{%
    \textbf{Definición.-} (Propiedad de sustitución bruta) Una función de exceso de demanda agregada $z(\cdot)$ satisface la propiedad de sustitución bruta si cuando $\mathbf p'$ y $\mathbf p$ sean vectores de precios tal que, para algún $\ell$, $p_\ell'<p_\ell$ y $p_k'=p_k$ para $k\neq\ell$, tenemos $z_k(\mathbf p')>z_k(\mathbf p)$ para $k\neq \ell$ tiene a lo sumo un equilibrio de intercambio; es decir, $z(p)=0$ tiene a lo sumo una solución.

Para motivar el concepto (y justificar su nombre), considérese la función de demanda de un consumidor en una situación con dos bienes. A precios dados, la matriz de sustitución de la demanda tiene entradas diagonales negativas y, como consecuencia, entradas fuera de la diagonal positivas: si se eleva el precio de un bien, la demanda compensada del otro bien aumenta. Sin embargo, si no descontamos los efectos riqueza (es decir, si observamos el efecto de los precios sobre la demanda no compensada), entonces es posible que un aumento en el precio de un bien reduzca la demanda de ambos bienes: en términos brutos, los dos bienes pueden ser complementarios. \textbf{¿Por qué esto supondría un problema para la unicidad global del equilibrio?}

Decimos que los dos bienes son sustitutos brutos si esto no ocurre, es decir, si un aumento en el precio de un bien reduce la demanda (no compensada o bruta) de ese bien y aumenta la demanda (no compensada o bruta) del otro bien. Por extensión, el mismo término se utiliza en el caso de $L$ bienes para la propiedad que afirma que cuando aumenta el precio de un bien, aumenta la demanda de todos los demás bienes (y, por lo tanto, disminuye la demanda de ese bien). Para L>2, sin embargo, esta no es en absoluto una propiedad necesaria ni siquiera de la demanda compensada. De hecho, la propiedad de sustitución bruta es muy restrictiva. No obstante, puede tener sentido para problemas con unos pocos bienes muy agregados o para aquellos en los que los bienes poseen simetrías especiales.
}

\textbf{Proposición.-} Una funciónes de exceso de demanda agregada $z(\cdot)$ que satisface la propiedad de sustitución bruta tiene como máximo un equilibrio de intercambio. Es decir, $z(\mathbf p)=0$ tiene como máximo una solución.

Para probarlo, basta con mostrar que $z(p)=z(p')$ no puede ocurrir cuando $p$ y $p'$ son dos vectores de precios que no son colineales. Por homogeneidad de grado cero, podemos suponer que $\mathbf p'\geq \mathbf p$ y que $p'_\ell=p_\ell$ para algún $\ell$. Ahora consideremos alterar el vector de precios $\mathbf p'$ para obtener el vector de precios $\mathbf p$ en $L-1$ pasos, reduciendo (o manteniendo inalterado) el precio de cada mercancía $k\neq \ell$ una a la vez. Por sustitución bruta, el exceso de demanda del bien $\ell$ no puede\ disminuir en ningún paso y, dado que $\mathbf p\neq\mathbf p'$, aumentará efectivamente en al menos un paso. Por lo tanto, $z_\ell(\mathbf p)>z_\ell(\mathbf p')$. 

En este escenario, gráficamente la FED agregada de una economía (reducida a dos bienes) se podría representar como:

\imagenfit{Unicidad_Global.png}{La Curva de Oferta de una FED que cumple con la propiedad de sustitución bruta: Unicidad global}{\textit{Microeconomis Theory}. MAS-COLELL et al. 1995}

Una característica importante de la propiedad de sustitución bruta es que es aditiva a través de las funciones de exceso de demanda. En particular, si las funciones individuales de exceso de demanda la satisfacen, entonces la función agregada también la satisface.

            \textcolor{red}{Podría esperarse establecer la unicidad en economías con producción aplicando la propiedad de sustitución bruta al exceso de demanda que incluya la producción. Sin embargo, el uso directo de la propiedad de sustitución bruta en un contexto de producción es limitado. Imagínese, por ejemplo, una situación en la que los insumos y los productos son bienes distintos: si el precio de un insumo aumenta, la demanda de todos los demás insumos puede disminuir, no aumentar como exigiría la propiedad de sustitución bruta, simplemente porque el nivel óptimo de producción disminuye. Indirectamente, sin embargo, el concepto de sustitutos brutos puede seguir siendo bastante útil.
            
            Entonces, suponiendo que el lado de la producción de la economía nos viene dado por una tecnología arbitraria $Y_j \subset \mathbb{R}^L$ del tipo de rendimientos constantes, convexa. ¿Qué condiciones que involucren únicamente el lado de la demanda de la economía garantizan la unicidad de las asignaciones de equilibrio? \textbf{A partir del análisis presentado en el PTFEB para un planificador benevolente, ya conocemos una respuesta: si una autoridad de bienestar se asegura de que la riqueza esté siempre distribuida de modo que maximice una función de bienestar social (estrictamente cóncava), entonces la economía admite un consumidor representativo normativo (estrictamente cóncavo), y el equilibrio corresponde necesariamente al único óptimo de Pareto de esta economía de un solo consumidor.} 
            
            Sin embargo, \textbf{en nuestro marco actual este no es un enfoque prometedor porque la riqueza se deriva de las dotaciones iniciales y solo por coincidencia podemos esperar que la distribución inducida de la riqueza maximice una función de bienestar social.} Por lo tanto, nos concentraremos en una \textbf{condición más débil¿?} y, para el propósito en cuestión, más interesante: que el \textbf{axioma débil de la preferencia revelada} se cumpla \textbf{para el exceso de demanda agregado} de los consumidores. Para ver todo el desarrollo y la demostración, se puede consultas [\authorfont{Ver Manual}].

            \textbf{Definición.-} (El axioma débil para funciones de exceso de demanda) La función de exceso de demanda $z(\cdot)$ satisface el axioma débil de la preferencia revelada (AD) si, para cualquier par de vectores de precios $\mathbf p$ y $\mathbf p'$, se cumple que
            $z(\mathbf p)\neq z(\mathbf p')$ y $\mathbf p\cdot z(\mathbf p')\leq 0$ implican $\mathbf p'\cdot z(\mathbf p)>0$.
            
            En palabras, la definición dice que si $\mathbf p$ es revelado preferido a $\mathbf p'$, entonces $\mathbf p'$ no puede ser revelado preferido a $\mathbf p$ [es decir, $z(\mathbf p)$ no puede ser asequible bajo $\mathbf p'$]. El axioma siempre es satisfecho por la función de exceso de demanda de un solo individuo, pero es una condición fuerte para el exceso de demanda agregado. ¿Qué hay de la suficiencia? El axioma débil no es una condición suficiente para la unicidad, pero sí garantiza que, para cualquier $Y$ convexo de rendimientos constantes, el conjunto de vectores de precios de equilibrio es convexo. Aunque esta propiedad de convexidad ciertamente no es lo mismo que unicidad, tiene una implicación inmediata de unicidad: si una economía tiene solo un número finito de equilibrios de precios (normalizados), el equilibrio debe ser único.
            
            ¿Qué condiciones sobre las preferencias y las dotaciones de los $I$ consumidores garantizan que la función de exceso de demanda agregada $z(\mathbf p)$ cumple el AD? Para comenzar con un caso relativamente simple, supongamos que todos los vectores de dotación $\omega_i$ son proporcionales entre sí; es decir, que $\omega_i=\alpha_i \overline \omega$, donde $\overline\omega$ es el vector de dotaciones totales y $\alpha_i\ge 0$ son participaciones con $\sum_i \alpha_i=1$. En tal economía, la distribución de la riqueza entre los consumidores es independiente de los precios. Normalizando los precios de modo que $\mathbf p\cdot \overline\omega=1$, la riqueza del consumidor $i$ es $\alpha_i$, y $z_i(\mathbf p)=x_i(\mathbf p,\alpha_i)-\omega_i$: si los niveles individuales de riqueza permanecen fijos, la satisfacción del AD por la demanda agregada (o el exceso de demanda), aunque restrictiva, no es implausible. Un supuesto de proporcionalidad sobre las dotaciones iniciales no es muy sostenible en un contexto de equilibrio general. Por lo tanto, es importante preguntar qué nuevos efectos están en juego cuando la distribución de dotaciones no satisface esta hipótesis. Desafortunadamente, resulta que la no proporcionalidad de las dotaciones puede reducir la probabilidad de que el exceso de demanda agregado satisfaga el axioma débil.
            
            Por último, ¿cuál es la relación entre la sustitución bruta y el axioma débil? Claramente, el axioma débil no implica la propiedad de sustitución bruta pues esta última puede violarse incluso en economías cuasilineales de un solo consumidor. La relación inversa no es tan obvia, pero no obstante es cierto que la propiedad de sustitución bruta no implica el axioma débil. Sin embargo, existe una conexión que es importante. La propiedad de sustitución bruta implica que, si $z(\mathbf p)=0$ y $z(\mathbf p')\neq 0$, entonces $\mathbf p\cdot z(\mathbf p')>0$
            } 

\subsubsection{Estabilidad: Precios}
Una característica distintiva que diferencia a la economía de otros campos científicos es que, para nosotros, las ecuaciones de equilibrio constituyen el centro de nuestra disciplina.\footnote{%
    Otras ciencias, como la física o incluso la ecología, ponen comparativamente mayor énfasis en la determinación de leyes dinámicas de cambio. En contraste, hasta ahora apenas hemos mencionado la dinámica. La razón, dicho informalmente, es que los economistas son buenos (o eso esperamos) para reconocer un estado de equilibrio, pero son pobres a la hora de predecir con precisión cómo evolucionará una economía en desequilibrio. Ciertamente existen principios dinámicos intuitivos: si la demanda es mayor que la oferta, entonces el precio aumentará; si el precio es mayor que el costo marginal, entonces la producción se expandirá; si los beneficios de la industria son positivos y no hay barreras a la entrada, entonces nuevas empresas entrarán, y así sucesivamente. La dificultad radica en traducir estos principios informales en leyes dinámicas precisas.
} La estática comparativa es la metodología analítica que se ocupa del estudio de cómo los equilibrios de un sistema se ven afectados por cambios en diversos parámetros del entorno. 

En un contexto general, el intento más famoso de esta traducción fue realizado por WALRAS (1874), y la versión moderna de sus ideas ha llegado a conocerse como la teoría de la estabilidad del tâtonnement. Revisamos brevemente dos modelos de tipo tâtonnement, uno de ajuste puro de precios y el otro de ajuste puro de cantidades. Debemos enfatizar, sin embargo, que estos son solo dos ejemplos. De hecho, una de las dificultades en este ámbito es la abundancia de modelos plausibles de desequilibrio. Aunque existe una única manera de estar en equilibrio, hay muchas maneras diferentes de estar en desequilibrio.

\paragraph{\textit{Tâtonnement} en precios: General y reforzado}
Siguiendo con el ejemplo de una economía de intercambio formalizada mediante una función de exceso de demanda $z(\cdot)$. Supongamos que tenemos un $\mathbf p$ inicial que no es un vector de precios de equilibrio, de modo que $z(p)\neq 0$. Por ejemplo, la economía puede haber sufrido un choque y $\mathbf p$ puede ser el vector de precios de equilibrio previo al choque. Definimos entonces los siguientes conceptos:
\begin{lnum}
    \item \textbf{Estabilidad local}. Llamamos a un equilibrio ($\mathbf p$) localmente estable si, siempre que el vector inicial de precios esté suficientemente cerca de él, la trayectoria dinámica hace que los precios relativos converjan a los precios relativos de equilibrio (el equilibrio es localmente totalmente inestable si cualquier perturbación conduce a que los precios relativos diverjan de los precios de equilibrio). Entonces, un equilibrio (regular) p_1/p_2 es localmente estable o localmente totalmente inestable según el signo de la pendiente del exceso de demanda en el equilibrio, es decir, según el índice del equilibrio (recuérdese la Definición 17.D.2). Si el exceso de demanda tiene pendiente negativa en p_1/p_2 (como en la Figura 17.H.1), entonces un ligero desplazamiento de p_1/p_2 por encima de p_1/p_2 generará exceso de oferta del bien 1 (y exceso de demanda del bien 2), y por lo tanto el precio relativo volverá a moverse hacia el nivel de equilibrio p_1/p_2. El efecto es el inverso si el exceso de demanda tiene pendiente positiva en p_1/p_2; y,
    \item \textbf{Estabilidad global}. Existe estabilidad del sistema si para cualquier posición inicial, $\mathbf p(0)$, la trayectoria correspondiente de los precios relativos $\mathbf p(t)$ converge a algún equilibrio de manera arbitrariamente cercana cuando $t\to\infty$.
\end{lnum}

Para probar su potencia, presentamos una versión en ecuaciones diferenciales propuesta por SAMUELSON (1947), que adopta la forma específica y de la que se extraen las siguientes conclusiones:

\eqblock{
\begin{aligned}
    \left.\begin{aligned}
        &\hspace{6em}\text{GENERAL}\\
        &\over\begin{aligned}
            \\
            &\hspace{2em}\underset{\substack{\\[0.5em]\\\text{\scriptsize
            $\text{\scriptsize donde,}\;\;
            \left\{\begin{aligned}
                &\mathrm dp_\ell/\mathrm dt&&\equiv \text{\scriptsize Ratio de ajuste para el bien $\ell$}\\[0.5em]
                &c_\ell>0&&\equiv\text{\scriptsize Cte. determinando la velocidad de ajuste}
            \end{aligned}\right.$}}}{\boxed{\frac{\mathrm dp_\ell}{\mathrm dt}=c_\ell z_\ell(\mathbf p)\qquad \forall \ell}}\\[1em]
            &\underset{\substack{\text{\scriptsize La única restricción}}}{\Longrightarrow}\hspace{1em}
            \substack{
            \text{Todas las trayectorias dinámicas}\\
            \text{quedan contenidas en la frontera}\\
            \text{definida por la propiedad (v) de la}\\
            \text{FDE agregada}}
        \end{aligned}
    \end{aligned}\hspace{1em}\right|
    \hspace{2em}\begin{aligned}
        &\hspace{6em}\text{PARTICULAR}\\
        &\over\begin{aligned}
            \\
            &\text{Si,}\qquad \{\text{Regular}\;\;\;\text{y}\;\;\;\;\text{AD}\;\;\;\;\text{y/o}\;\;\;\;\text{Sustitución bruta}\}\quad (\checkmark)\\[2em]
            &\text{Entonces,}\quad  \exists! \mathbf p^*\in \mathbb R^L: z(\mathbf p^*)=0\\[0.5em]
            &\Longrightarrow\forall \mathbf p\in \mathbb R^L: \;\mathbf p\neq\mathbf p^*\;\;\wedge\;\;\mathbf p\neq\alpha\mathbf p^*,\;\;\alpha\in\mathbb R,\\[2em]
            &\hspace{6em}\boxed{\lim_{t\to\infty}\mathbf p(t)=\mathbf p^*,\;\;\;\forall \mathbf p}
        \end{aligned}
    \end{aligned}
\end{aligned}
}{\textit{Tâtonnement} en precios: Caso particular y general. Propiedades positivas: Estabilidad}

Como vemos, por simple que se la intuición económica que subyace en el proceso de estabilidad o convergencia al equilibrio: por ejemplo, que si la dotación de un bien aumenta, entonces su precio de equilibrio disminuye, o que atendiendo al principio generalmete aceptado de oferta y demanda que sugiere que los precios se ajustarán al alza para los bienes con exceso de demanda y a la baja para aquellos con exceso de oferta. En un contexto de equilibrio general, se requieren condiciones fuertes para que estos efectos se cumplan. A estas alturas esto no debería sorprendernos: ya sabemos que los efectos riqueza y/o la falta de suficiente sustituibilidad pueden socavar efectos intuitivos de estática comparativa.\footnote{%
    De hecho, existen otras cuestiones dificilmente justificables en esta versión diferencial que no entraremos a vlorar pero si cabe discutir.

Entre otras, ¿qué agente económico está a cargo de los precios? Más aún, ¿por qué debe cumplirse la ley del precio único fuera del equilibrio (es decir, por qué bienes idénticos deben tener precios idénticos fuera del equilibrio)? ¿Qué tipo de tiempo representa $t$? No puede ser tiempo real, porque, tal como está el modelo, un $\mathbf p$ de desequilibrio no es compatible con la factibilidad (es decir, no todos los planes de consumo pueden realizarse simultáneamente).

Quizá la respuesta más sensata a todas estas preguntas es que debe pensarse no como un modelo de la evolución real de una economía impulsada por la oferta y la demanda, sino más bien como un proceso tentativo de prueba y error que tiene lugar en un tiempo ficticio y es dirigido por un agente de mercado abstracto empeñado en encontrar el nivel de precios de equilibrio (o, más modestamente, empeñado en restablecer el equilibrio tras una perturbación). Este es de hecho el significado original que se le dió: \textit{subastador walrasiano}.
} 

\paragraph{\textit{Tâtonnement} en cantidades: General y reforzado}
En el análisis hasta ahora, los precios podían estar fuera de equilibrio, pero las cantidades, es decir, las cantidades demandadas y ofrecidas, están siempre en sus valores de equilibrio (esto es, de maximización de utilidad y de beneficios en caso de incluir producción). No obstante, el mismo proceso de convergencia puede analizarse con respecto a cantidades, considerando que son las cantidades, y no los precios, las que pueden estar en desequilibrio. 

Esto exige reformular el modelo para mostrarlo de manera rigurosa. En particular, reformularlo de manera tal que se incluya el proceso de transformación/producción explícitamente. Las conclusiones, no obstante, serán las mismas que las anteriores.\footnote{%
    De manera simplificada, supongamos que existe un único conjunto de producción Y. (No hay ningún problema con considerar varios, es solo por simplicidad, incluso puede considerarse como el conjunto de producción de toda la economía). En cualquier momento del tiempo, suponemos que está dado un único vector de producción fijo $y\in Y$. Los precios, sin embargo, están siempre en equilibrio en el sentido de que el sistema de equilibrio general de la economía, condicionado a $y$, genera algún sistema de precios de equilibrio $\mathbf p(y)$ (es decir, procedemos en el corto plazo como si el conjunto de producción de corto plazo fuera $\{y\}-\mathbb R^L_+$). Esto describe el equilibrio de corto plazo de la economía.

¿Cuál es una dinámica apropiada para esta economía? Tiene sentido pensar que, sea cual sea, el cambio en la producción en el tiempo $t$, $\mathrm dy(t)/\mathrm dt\in \mathbb{R}^L$, mueve la producción en una dirección que incrementa los beneficios cuando el vector de precios en el tiempo $t$, $\mathbf p(y(t))$, se toma como dado:

\textbf{Definición.-} Decimos que la trayectoria diferenciable $y(t)\in Y$ es admisible si $\mathbf p(y(t))\cdot (\mathrm dy(t)/\mathrm dt)\ge 0$ para todo $t$, con igualdad solo si $y(t)$ maximiza beneficios para $\mathbf p(y(t))$ (en cuyo caso podríamos decir que estamos en un equilibrio de largo plazo).

Una diferencia entre los enfoques de tâtonnement de precios y de cantidades que añade atractivo al segundo es que la factibilidad queda ahora asegurada en todo $t$ y que, como resultado, podemos interpretar la dinámica como ocurriendo en tiempo real.(No obstante, es importante darse cuenta de que, incluso entonces, este no es un modelo plenamente dinámico: los problemas de optimización de los consumidores siguen siendo estáticos y libres de retroalimentaciones expectacionales y las empresas siguen reglas ingenuas de ajuste de corto plazo (en un espíritu más positivo, podría llamarse a esto comportamiento adaptativo, más que ingenuo)). ¿Llevará necesariamente una trayectoria admisible al equilibrio de largo plazo? No podemos explorar realmente esta cuestión ya que, como es habitual, la respuesta es solo bajo circunstancias especiales. Un ejemplo limitado, pero importante resulta de que: Si existe un único consumidor estrictamente convexo, entonces cualquier trayectoria admisible converge al (único) equilibrio. Para ver la demostración consultese [\auhorfont{Ver Manual}]. De forma simplificada, dado que la utilidad es creciente, debemos necesariamente alcanzar el vector de producción $y$ en el que la utilidad se maximiza en el conjunto factible de producción (es decir, el equilibrio).
} Por ello, resulta mucho más adecuada hablar directamente de las necesidad o razones prácticas que motivarían el análisis de la estabilidad desde un desequilibrio en cantidades y no precios: ¿por qué es importante introducir estos mecanismos de ajuste diferentes, teniendo además en cuenta que la mayoría de ajustes se realizan mediante precios, pues son las empresas quien ofertan sus productos a un precio determinado?

Bajo las condiciones características que garantizan la existencia de un equilibrio único y global, podemos ilustrar con un ejemplo sencillo la importancia de disponer de estas dos herramientas analíticas:


\imagenfit{CasosPolares_AjustesDiferentes.png}{Dos casos polares con elasticidad cantidad-precio diferentes: Oferta y Demanda}{Apuntes ICEX-CECO. Tema 3.A.21}

¿Sería adecuado en ambos escenarios utilizar el mismo proceso de ajuste? ¿Qué resultados obtendríamos ajustando vía precios en el primer mercado (arriba)? ¿Y en el segundo (abajo)? 

En el primer mercado, suponiendo que existe un exceso de demanda positivo, sabemos que los precios deben disminuir para volver al equilibrio. 
 
En este caso, el ajuste adecuado sería por cantidades, puesto que la relación entre la función de exceso de demanda y el precio no es inversa, sino lineal. 
 
En este otro caso, podremos utilizar un mecanismo de ajuste vía precios puesto que se cumple la condición de relación inversa entre la función de exceso de demanda y el nivel de precios. 
Algunos ejemplos de éste tipo de mercados se encuentran fundamentados empíricamente, como son:
	Los mercados en los que se observa target earning. Cuando los precios varían (con frecuencia o no) por shocks transitorios exógenos al mercado, podríamos observar comportamientos “estratégicos” de la ofertantes, de tal forma que exista una curva de oferta inversa, donde a menor precio, mayor cantidad ofertada (oferta extensiva) y a mayor precio menor cantidad (oferta intensiva, por ganancias ya capitalizadas).  
	Mercados de commodities con productores con objetivos/restricciones de ingreso.  Algunos mercados, como el del petróleo (en algunos casos) presentan cóncavidades hacia el origen a partir de una cierta cantidad de producción. Dichos mercados presentan la característica de una demanda muy inélastica en comparación con la oferta. Por tanto, la mayoría de ajustes se producen vía cantidades y no a través de precios. 
	Los mercados mayoristas de electricidad, son un buen ejemplo de adaptación contínua mediante precios. Se caracterizan por tener precios marginales locacionales (LMP, por sus siglas en inglés)    
Por tanto, como se observa, para demostrar el cumplimiento de tres propiedades positivas de relevancia como son la existencia, unicidad y estabilidad del equilibrio general competitivo se ha de recurrir a supuestos muy restrictivos y deben adaptarse adecuadamente. 

\subsection{Implicaciones de Política Económica: Enfoque normativo}
Tan importante resulta comprender cómo se obtiene el equilibrio general o el cumplimiento de ciertas propiedades positivas, como pasar a analizar desde un punto de vista normativo la deseabilidad del equilibrio general.  Así que, ¿qué puede decirse acerca de las propiedades normativas de dicho equilibrio general?
La teoría del bienestar permite relacionar los conceptos de eficiencia de Pareto y equilibrio general competitivo a través de los Teoremas Fundamentales de la Economía del Bienestar.
2.4.1. Primer Teorema Fundamental de la Economía del Bienestar
El primer teorema establece que, si un mercado funciona de forma competitiva, el equilibrio general obtenido será Pareto eficiente (que no equitativo). Los precios dan los incentivos para que los agentes lleguen a un resultado eficiente actuando en su interés propio. 
Dicho de otra manera, el vector de precios del equilibrio general competitivo garantiza no solamente que todos los agentes estén llevando un comportamiento optimizador, sino que, además, no es posible mejorar a ninguno de ellos sin que al menos otro empeore.
Este primer teorema se cumplirá siempre que las preferencias de los consumidores muestren no saciabilidad local. (no está esto en contradicción con la condición de unicidad del equilibrio?)
2.4.2. Segundo Teorema Fundamental de la Economía del Bienestar
Establece que toda asignación eficiente en sentido de Pareto es alcanzable a partir del equilibrio general a través de una redistribución previa de los recursos mediante impuestos y subvenciones de suma fija.
Este teorema tiene importantes implicaciones desde el punto de vista de política económica, pues permite conciliar los criterios de equidad y eficiencia. Es decir, incluso en ausencia de fallos de mercado, el sector público tiene margen de actuación para alcanzar asignaciones eficientes en sentido de Pareto que se consideren socialmente más deseables. Para ello es necesario redistribuir previamente las dotaciones iniciales mediante impuestos y subvenciones de suma fija y dejar actuar libremente a los mecanismos del mercado. El cumplimiento de este segundo teorema requiere la convexidad de la tecnología de producción y de las funciones de utilidad de los consumidores.
2.4.3. Implicaciones
A pesar de las importantes implicaciones en materia de política económica de los teoremas fundamentales de la economía del bienestar, es evidente que estos presentan limitaciones que les restan validez práctica:
	Requieren que se cumplan todos los supuestos competitivos, pero en la práctica existen fallos de mercado que invalidan las conclusiones aquí presentadas.
	Incluso en un contexto en el que se cumplen todos los supuestos competitivos, la intervención del policymaker con base en el segundo teorema fundamental requeriría de mucha información para poder llevar a cabo la redistribución inicial de recursos. En concreto, se necesitaría ser plenamente conocedor de las dotaciones iniciales de todos los agentes y de cuáles deberían ser los impuestos y subvenciones a aplicar sobre ellos para que el mecanismo de mercado permita llegar a la asignación deseada.
	Incluso disponiendo de toda la información, la intervención debe hacerse mediante impuestos y subvenciones de suma fija, figuras no utilizadas en la práctica por sus limitaciones inherentes.


\clearpage
\section{Equilibrio General: Fallos de Mercado}
\puntosclave{
    En el enfoque del núcleo, los agentes negocian directamente entre ellos, teniendo un rol diferente al de la obra de Walras. Cuando el número de agentes tiende a infinito, ambos enfoques llevan a la misma asignación de equilibrio.
    
	Arrow y the Broad desarrollan el concepto de bienes contingentes para introducir incertidumbre en el modelo de equilibrio general. La posibilidad de comerciar bienes contingentes lleva en el equilibrio a un reparto eficiente del riesgo.
    
	En el enfoque de patinkin, que introduce dinero en el modelo de equilibrio general, cambios en la cantidad de dinero circulando en la economía no tienen efectos reales. Estos cambios se limitan a cambiar el precio absoluto de los bienes en una proporción idéntica para todos ellos.
    
	Para que los resultados y corolarios del modelo de equilibrio general sean completamente independientes del número de consumidores, es crucial el cumplimiento del supuesto de convexidad de preferencias.
}
	
Tras todo lo expuesto, resulta fácil entender la potencia de la Teoría del Equilibrio General. La obra que inició WALRAS demuestra que, si se cumplen un número muy reducido de condiciones, existe un vector de precios tal que se alcanza un equilibrio óptimo en el sentido de PARETO (máxima eficiencia), garantizando el vaciamiento de los mercados. Además, proporciona una herramienta analítica muy potente para reducir la complejidad analítica inicial a una Función de Exceso de Demanda que permite determinar el vector de precios relativos a un bien, el \textit{numerario}.

No obstante, el contexto económico en el que opera es una versión idealizada de la realidad. Existen ciertas restricciones fundamentales, tanto institucionales como interactivas, que \textit{a priori} limitan la aplicabilidad de este marco teórico. Por lo tanto, en esta segunda parte analizaremos brevemente otros enfoques que permiten dar \textbf{aproximar la teoría a la realidad} de diferentes formas. ¿Es esto posible? Por ejemplo, ¿qué pasa cuando la información no es completa? ¿qué pasa cuando los agentes no conocen todos los precios? ¿estos sesgos no desvian su conducta de lo que sería óptimo para ellos? Más aún, ¿qué pasa cuando los mercados no son completos? O, ¿y si las empresas ostentan poder de mercado?

Todas estas preguntas se han ido estudiando progresivamente por la ciencia económica utilizando instrumentos de diversos campos.

\subsection{Incertidumbre: Información incompleta}
Uno de los supuestos más alejados de la realidad es considerar que los agentes disponen de información completa sobre todos los fundamentos de la economía en su conjunto. Es evidente que muchos fundamentales son inciertos, de hecho, bien podría ocurrir que los elementos constitutivos de la economía fuesen en sí mismos inciertos: las tecnologías, las dotaciones y las preferencias. Sobra decir la evolución dinámica de esta economía por parte de los agentes. La incertidumbre ha sido tratado en un marco general por diversos autores. 

\subsubsection{Información simétrica: Mercado de activos}
El marco de ARROW-DEBREU fue uno de los primeros en introducir la incertidumbre y proporciona gran ejemplo del poder de la teoría del equilibrio general.\footnote{%
    En esencia, el equilibrio de ARROW DEBREU consiste en lo siguiente. 

Supongamos la existencia de un mercado para cada bien contingente $\ell_s$. Es decir, un mercado para cada bien para cada estado de la naturaleza posible de la economía. Estos mercados se abren antes de la resolución de la incertidumbre, en la fecha $0$, podríamos decir. El precio del bien se denota por $p_{\ell_s}$. Lo que se compra (o se vende) en el mercado del bien contingente $\ell_s$ son compromisos de recibir (o de entregar) cantidades del bien físico $\ell$ si, y cuando, ocurre el estado del mundo $s$. Por tanto, obsérvese que, aunque las entregas son contingentes, los pagos no lo son. Nótese también que, para que este mercado esté bien definido, es indispensable que todos los agentes económicos sean capaces de reconocer la ocurrencia de $s$. Es decir, la información debe ser simétrica entre los agentes económicos. Formalmente, la economía de mercado que acabamos de describir no es más que un caso particular de las economías que hemos estudiado en la primera parte de la exposición. Por lo tanto, podemos aplicar a nuestra economía de mercado el concepto de equilibrio walrasiano y, con él, toda la teoría desarrollada hasta ahora: 
$$\begin{aligned}
    &\text{Equilibrio de ARROW-DEBREU}\\[1em]
    &\text{Se el conjunto,}\quad (x_1^*,\ldots,x_I^*,y_1^*\ldots,y_J^*)\in X_1\times\cdots\times X_I\times Y_1\times \cdots \times Y_J\subset\mathbb R^{LS(I+J)}\\[0.5em]
    &\text{y un sistema de precios contingente}:\mathbf p=(p_{11},\ldots,p_{LS})\in\mathbb R^{LS}\\[0.5em]
    &\mathbf p\equiv\text{Equilibrio ARROW-DEBREU}\iff
    \left\{\begin{aligned}
        &\text{(1)}\;\;\forall j,\;y_j^* \Longrightarrow\mathbf p\cdot y_j^*\geq \mathbf p\cdot y_j,\;\;\forall y_j\in Y_j\\[0.5em]
        &\text{(2)}\;\;\forall i,\;x_i^* \;\;\text{es máximo dado el conjunto presupuestario}\\[0.5em]
        &\text{(3)}\;\;\sum_ix_i^*=\sum_jy_j^*+\sum_i\omega_i
    \end{aligned}\right.
\end{aligned}$$
Al tratar con bienes contingentes es habitual denominar equilibrio de ARROW-DEBREU al equilibrio walrasiano. Esta nueva definición de equilibrio, que en esencia solo incorpora los estados contingentes en la asignación de equilibrio y el vector de precios, cumple con todas las implicaciones tanto positivas como normativas presentadas en la primera parte. No obstante, existe alguna diferencia en relación a los supuestos base: (i) especialmente, en el contexto presente, el supuesto de convexidad adquiere una interpretación en términos de aversión al riesgo (por ejemplo, en el marco de utilidad esperada la relación de preferencias $\succeq_i$ es convexa si las utilidades de BERNOULLI $u_{is}(x_{is})$ son cóncavas; y, (ii) para cualquier plan de producción, el beneficio de una empresa, $\mathbf p\cdot y_j$, es una cantidad no aleatoria de dinero. Es decir, las producciones y las entregas de bienes sí dependen, por supuesto, del estado del mundo, pero la empresa participa en todos los mercados contingentes y logra, por así decirlo, asegurarse completamente (esto tiene implicaciones importantes para la justificación de la maximización de beneficios como objetivo de la empresa).

} Representando la incertidumbre como un \textbf{conjunto exhaustivo de estados del mundo}, $(s_1,\ldots,s_S)\in\mathbb R^S$. La incertidumbre emerge de la probabilidad de que se de uno u otro estado de la naturaleza, sin que se conozca con certeza cuál va a ser. 

En este escenario, el equilibrio ARROW-DEBREU comparte exactamanete las mismas propiedades que el expuesto en nuestra primera parte. No obstante, la condición paretiana tiene unas implicaciones ligeramente diferentes. No se trata ahora de la asignación eficiente de todos los recursos de la economía, sino de la asignación eficiente de todas los escenarios/estados contingentes posibles. Es decir, se trata ahora de una \textbf{asignación eficiente del riesgo}.\footnote{%
    El riesgo como incertidumbre respecto a los estados contingentes posibles. Al ser los mercados contingentes completos (un mercado para cada estado y bien comerciable), la información simétrica (todos preveen el mismo escanario en cada estado) y cumplir las condiciones iniciales expuestas en la primera parte relativas a los supuestos de partida (institucionales, fundamentales y de comportamiento), entonces la asignación del riesgo/incertidumbre es eficiente.
} No obstante, fue criticado por partir de un supuesto difícilmente realista: todo el comercio tiene lugar simultáneamente y antes de que se resuelva la incertidumbre; es, por así decirlo, un asunto de una sola vez. 

En realidad, el comercio tiene lugar de forma secuencial a lo largo del tiempo, y con frecuencia como consecuencia de revelaciones de información. Las aportaciones de RADNER permiten reconciliar este supuesto con el equilibrio de ARROW-DEBREU. Para ello, define qué significa el comercio secuencial y demuestra las condiciones de equivalencia; es decir, bajo qué condiciones se alcanzan las asignaciones de equilibrio ARROW-DEBREU a través del comercio secuencial. 

\paragraph{Equilibrio de RADNER}
En primer lugar, una condición \textit{sine qua non} para que exista comercio secuencial es que este sea deseable, de manera que la asignación alcanzable en un contexto estático no pueda ser eficiente en el sentido de PARETO. Una causa clara de este escenario sería que el comercio ex post puede solventar la asuencia del mercado ex ante para algún bien, entonces es posible que mediante comercio secuencial se alcance una asignación eficiente. Para que esto ocurra, RADNER impone las siguientes condiciones:

\eqblock{
\begin{aligned}
    &\text{Condiciones de RADNER}\\[0.5em]
    &\underbrace{\left(\begin{aligned}
        &\text{(1)}\;\;\exists \mathbf z=(z_1,\ldots,z_S)\in\mathbb R^S&&\text{mercado contingente en $t=0$}\\[0.5em]
        &\text{(2)}\;\;\exists \mathbf x=(x_{11},\ldots,x_{1L}, \ldots,x_{S1},\ldots,x_{SL})\in\mathbb R^{LS}&&\underset{\text{\scriptsize OJO, la contingencia ya se ha resuelto en $t=0$, aquí es mercados dado el estado}}{\text{mercados spot en $t=1$, dado el estado}}\\[0.5em]
        &\text{(3)}\;\;\text{HER precios en $t=1$}&&\text{previsión completa: información simétrica}
    \end{aligned}\right)}\\[1em]
    &\hspace{20.5em}\Downarrow\\[1em]
    &\hspace{14em}\begin{aligned}
        &\underset{\text{\scriptsize vector de precios para el bien (activo) contingente comerciado en $t=0$}}{\exists \mathbf q=(q_1,\ldots,q_S)\in\mathbb R^S,\;\;\text{para}\;\;\mathbf z}\\[0.5em]
        &\underset{\text{\scriptsize vector de precios esperado para mercado spot del resto de mercados, en $t=1$}}{\exists \mathbf p=(\mathbf p_1,\ldots,\mathbf p_S)\in\mathbb R^{LS},\;\;\text{para}\;\;\mathbf x}
    \end{aligned}\\[1em]
    &\hspace{20.5em}\Downarrow\\[1em]
    &\hspace{5em}\begin{aligned}
        \forall i,\;\;\exists
        \left\{\begin{aligned}
            &(\mathbf x_{1i},\ldots,\mathbf x_{Si})\in\mathbb R^{LS}\\
            &(z_{1i},\ldots,z_{Si})\in\mathbb R^{S}
        \end{aligned}\right\},\;\;
        \boxed{
        \begin{aligned}
            &\max\;U_i(\mathbf x_{1i},\ldots,\mathbf x_{Si})\\[0.5em]
            &\text{s.a.}\;\;
            \left\{\begin{aligned}
                &\text{(1)}\;\;\sum_sq_sz_{si}\leq 0\\
                &\text{(2)}\;\;\mathbf p_s\cdot \mathbf x_{si}\leq \mathbf p_s\cdot \omega_{si}+\mathbf p_{1s}\cdot z_{si},\;\;\forall s
            \end{aligned}\right.
        \end{aligned}}\;\;(\checkmark)  
    \end{aligned} 
\end{aligned}
}{Condiciones de RADNER: Exigimos HER (autocumplidas) para cada $s$. Teoría del Equilibrio General: Incertidumbre}

Por lo tanto, verificamos que si: (i) al menos un bien físico pueda comerciarse de manera contingente en $t=0$; (ii) todos los mercados spot existen en $t=1$, dado cualquier estado; y, (iii) los precios de equilibrio al contado son anticipados correctamente en $t=0$.\footnote{%
    Las expectativas de los consumidores deben ser autocumplidas, o racionales; es decir, exigimos que las expectativas de los consumidores acerca de los precios que despejarán los mercados al contado para los distintos estados $s$ efectivamente los despejen una vez que llega la fecha $t=1$ y se revela un estado $s$.
} Entonces, existe un vector de precios del bien contingente y un vector de precios en los mercados spotfrente a precios, que el problema del consumidor se puede plantear frente a una dotación de cartera y de consumo en un número de mercados a plazo requeridos que pasa de $LS$ (ARROW-DEBREU) a $S$ (RADNER).\footnote{%
    Enfrentado a un vector de precios del bien contindente $\mathbf q\in\mathbb{R}^S$ en $t=0$ y a los precios al contado esperados $(\mathbf p^1,\ldots,\mathbf p^S)\in\mathbb{R}^{LS}$ en $t=1$, cada consumidor $i$ formula un plan de consumo, o de comercio, $(z_{1i},\ldots,z_{Si})\in\mathbb{R}^S$ para los bienes contingentes en $t=0$, así como un conjunto de planes de consumo en los mercados al contado $(\mathbf x_{1i},\ldots,\mathbf x_{Si})\in\mathbb{R}^{LS}$ para los distintos estados que pueden ocurrir en $t=1$. La intuición de este resultado es razonablemente directa: si el comercio al contado puede tener lugar dentro de cada estado, entonces la única tarea que resta en $t=0$ es transferir de forma eficiente el poder adquisitivo total del consumidor entre estados. Como esto puede lograrse utilizando comercio contingente en un solo bien, somos capaces de reducir el .
} 

Por último, si el conjunto de asignaciones de cartera y de consumo que resuelven el problema de maixmización cumplen a su vez la siguiente condición:

\eqblock{
\begin{aligned}
    &\text{s.a.}\quad \sum_iz_{si}^*\leq0\quad \text{y,}\quad \sum_ix_{si}^*\leq \sum_i\omega_{si}\;\;\forall, i\\[1em]
    &\text{Eq RADNER}\Longrightarrow\text{Eq. ARROW-DEBREU}\quad \blacksquare
\end{aligned}
}{Condiciones de RADNER: Equivalencia RADNER y ARROW-DEBREU. Teoría del Equilibrio General: Incertidumbre}

Entonces se cumple que el equilibrio de RADNER equivale al equilibrio de ARROW-DEBREU y, en consecuencia, además de todas las implicaciones positivas expuestas en la primera parte, la asignación del riesgo es óptima.

\paragraph{Desarrollo formal}
Ahora bien, el constructo coneptual del activo subyacente en el modelo anterior permite la transferencia intertemporal de riqueza de manera independiente entre estados: podemos escoger, libremente, en cómo distribuimos nuestro riesgo entre estados, a voluntad. Podría decirse, entonces, que constituye un \textit{modelo idealizado} del mercado de futuros, donde cada agente distribuye su riesgo no solo atendiendo a sus preferencias sino siendo perfectamente distribuido. Es decir, hay un contrato distinto para cada estado que determina nuestra riqueza en términos del bien contingente si ocurre dicho estado.\footnote{%
    Eso significa que, en $t=0$, cada agente decide con certitud cuánto riesgo/riqueza quiere en cada estado por separado.
} 

En la realidad, no obstante, esto no ocurre así. En la realidad existen activos que nos permiten \textbf{distribuir} el riesgo conforme a nuestros intereses. Es decir, nos enfrentamos a problemas económicos iguales pero mediante \textbf{instrumentos imperfectos}, que ya no permiten replicar perfectamente nuestras preferencias, sino distribuirlas en el espacio.

En este escenario, lo máximo a lo que podemos aspirar es a distribuir el riesgo mediante activos cuyos rendimientos vienen condicionados al estado de la naturaleza que realmente ocurra. De esta forma, sea $r^k=(r^k_1,\ldots,r_S^k)\in\mathbb R^S$ el vector de rendimientos del activo $k$, y $K$ el conjunto/estructura de activos. Las condiciones de equilibrio de RADNER se modifican de forma que:

\eqblock{
\begin{aligned}
    &\hspace{8em}\text{Condiciones de RADNER (pero, variaciones siguiente)}\\[1em]
    &\left\{\begin{aligned}
        &\text{(1)}\;\;\exists\mathbf z^*_i=(z^*_{1i},\ldots,z^*_{Ki})\in\mathbb R^K&&\underset{\text{\scriptsize teniendo en cuenta que $\forall k$,  $r^k=(r_1^k,\ldots,r_S^k)\in\mathbb R^S$}}{\text{portfolios óptimos $\forall i$}}\\[1em]
        &\text{(2)}\;\;\exists\mathbf q=(q_1,\ldots,q_K)\in\mathbb R^K&&\text{precios de trade en $t=0$}\\[1em]
        &\text{(3)}\;\;\text{$R=S\times K$ es completa, i.e. rank $R=S$}&&\begin{aligned}
            \text{matriz de retornos de la estructura}\\
            \text{de activos y estados de naturaleza}
        \end{aligned}
    \end{aligned}\right\}\\[1em]
    &\hspace{5em}\Downarrow\text{Eq. RADNER $=$ Eq. ARROW-DEBREU}\;(\checkmark)\\[1em]
    &\begin{aligned}
        \boxed{\begin{aligned}
            &\max\;U_i(\mathbf x_{1i},\ldots,\mathbf x_{Si})\\[0.5em]
            &\text{s.a.}\;\;
            \left\{\begin{aligned}
                &\text{(1)}\;\;\sum_kq_kz_{ki}\leq 0\\
                &\text{(2)}\;\;\mathbf p_s\cdot \mathbf x_{si}\leq \mathbf p_s\cdot \omega_{si}+\sum_{k}\mathbf p_{1s}\cdot z_{ki}\cdot r_{sk}\;,\;\;\forall s
            \end{aligned}\right.
        \end{aligned}}\quad+\quad 
        \boxed{\sum_{i}z_{ki}^*\leq 0\;\;\text{y,}\;\;\sum_ix_{si}^*\leq\sum_i\omega_i,\;\;\;\forall i}
    \end{aligned}
\end{aligned}
}{Condiciones de RADNER: Mercado de activos. Teoría del Equilibrio General: Incertidumbre}

Por tanto, vemos como con la introducción de activos (como instrumentos de transferencia del riesgo), la importancia del concepto de completitud es total. Con una estructura de activos completa, los agentes económicos están, en efecto, libres de restricciones en sus transferencias de riqueza entre estados (salvo, por supuesto, las impuestas por sus restricciones presupuestarias). Por lo tanto, en equilibrio, sus elecciones de cartera inducen los mismos consumos en el segundo período que en el equilibrio de Arrow–Debreu, y así se alcanza la plena optimalidad de Pareto.

\subsubsection{Implicaciones de Política Económica}
Este resultado tiene implicaciones fortísimas ya que permite generalizar las propiedades positivas y normativas de la Teoría presentada en la primera parte del tema a una economía mucho más compleja, con multitud de periodos y con transferencias de riqueza intertemporales que distribuyen el riesgo (incertidumbre) entre estados por cada activo y, por tanto, no constituyen per se vehículos perfectos.

\paragraph{Mercados Incompletos: RADNER $\not\to$ ARROW-DEBREU (aleatorio)}
No obstante, exige que impongamos una restricción para que pueda existir un equilibrio ARROW-DEBREU: los mercados deben ser completos. En la realidad, los mercados pueden no existir por diversas razones, entre otras:
\begin{lnum}
    \item \textbf{Información asimétrica}. Una clase de razones se refiere a las asimetrías de información: los contratos para la entrega de bienes solo pueden hacerse contingentes a estados cuya ocurrencia pueda verificarse de manera satisfactoria para todas las partes contratantes;
    \item \textbf{Costes de transacción}. Otra clase de razones proviene de los costos de transacción: la disponibilidad de un mercado es, al fin y al cabo, de la naturaleza de un bien público; y,
    \item \textbf{Restricciones de exigibilidad}. Otra variedad de razones surge de las restricciones de exigibilidad: una promesa de entregar una unidad de un bien es inútil si la entrega no puede hacerse cumplir.\footnote{%
        Un ejemplo de una situación en la que la exigibilidad se vería facilitada es cuando tratamos con acciones de una empresa (el total de dotaciones de este activo es, por lo tanto, positivo) y no son posibles las ventas en corto. En ese caso, la exigibilidad está garantizada porque las acciones físicas (los derechos de propiedad sobre la empresa) se intercambian efectivamente en $t=0$. Este tipo de consideración, que lamentablemente entra en conflicto con el supuesto tan conveniente de que las ventas en corto ilimitadas son posibles, podría explicar por qué algunos activos pueden existir y otros no: para que un activo exista, ayuda que la variable aleatoria tenga una realidad física que pueda intercambiarse en $t=0$.
    }
\end{lnum}

Por tanto, cuando la estructura de activos no es completa atendiendo a los posibles estados de la naturaleza, $K < S$: un equilibrio de RADNER no tiene por qué ser óptimo de Pareto.\footnote{%
    Esto no es sorprendente: si las posibilidades de transferir riqueza entre estados son limitadas, entonces bien puede haber una pérdida de bienestar debido a la incapacidad de diversificar los riesgos en la medida que sería conveniente. Basta considerar el caso extremo en el que no hay activos en absoluto. Es decir, la equivalencia se rompe y no podemos garantizar la optimalidad paretiana en este contexto, será cuestión de azar.
}

\paragraph{Mercados incompletos: El papel de la empresa}
En esta segunda parte de la exposición nos hemos concentrado en el estudio de economías de intercambio y, esta vez, no ha sido solo por simplicidad. Si consideramos el mismo escenario presentado anteriormente: (i) comercio secuencial; y, (ii) vehiculos de transferencia son activos con rendimiento. 

Consideremos lo siguiente: (i) un entorno con dos períodos, $t=0,1$; (ii) dos estados posibles de la naturaleza en $t=1$, ($s_1,s_2$); (iii) introducimos la producción y consideramos una economía descentralizada en la que existen participaciones $\theta_i\ge 0$, con $\sum_i \theta_i=1$ sobre rendimientos de la empresa (nótese que los rendimientos son un activo que sirve como vehiculo de transferencia del riesgo: $a=(a_1,a_2)$); y, (iv) adoptamos el enfoque natural y consideramos que la producción contingente (dependiente del estado) de la empresa equivale, precisamente, a sus rendimientos. ¿Cómo podemos formular el problema de la empresa?

A priori, sabemos que el conjunto $A\subset\mathbb{R}^2$ de posibles elecciones de vectores de producción $(a_1,a_2)\in A$ de la empresa está acotado por rango. Gráficamente:

\imagenfit{Conjunto_A_empresa.png}{Ejemplo de posibilidades de producción para la empresa}{\textit{Microeconomic Theory}. MAS-COLELL et all, 1995}

Además, el vector de producción/rendimientos $(a_1,a_2)$ se elige antes de que se abran los mercados financieros del período $t=0$ por lo que la decisión es tomada por los propietarios iniciales.\footnote{%
    Dado que las acciones pueden venderse en $t=0$, los propietarios al final de $t=0$ pueden ser un conjunto distinto.
} Por tanto, ¿qué plan de producción elegirán los propietarios iniciales?

Resulta que la respuesta tiene fuertes implicaciones y depende crucialmente de la relación que exista entre su vector de rendimientos y el vector de rendimientos total $R$:\footnote{%
    \textbf{IMPORTANTE.-} Recordar que la matriz $R$ recoge los rendimientos contingentes de todos los activos que ya existen en la economía: cada columna es un activo; cada fila, un estado. Es decir, $R$ describe qué perfiles de riqueza futura pueden construirse combinando activos existentes.
}
\begin{la}
    \item \textbf{Si $A\subset \text{rango}(R)$}. Si el activo pertenece al rango de la matriz de rendimientos de la economía, entonces el perfil de rendimientos que genera la empresa ya puede replicarse exactamente mediante una cartera de activos existentes. Es decir, existe una combinación de activos tal que paga lo mismo que la empresa en cada estado. Por tanto, la empresa no introduce ningún riesgo nuevo en la economía y el resultado del problema de producción se vuelve tautológico: todo consumidor-propietario preferirá el plan de producción/rendimientos con mayor valor de mercado, \textbf{la empresa se vacía de contenido económico}; y,
    \item \textbf{Si $A\not\subset \text{rango}(R)$}. Si el activo no pertenece al rango de la matriz de rendimientos de la economía, entonces la empresa genera un perfil de rendimientos que no puede reproducirse con los activos existentes. Es decir, la empresa introduce un nuevo tipo de riesgo, una nueva dimensión de redistribución de riqueza entre estados. Por tanto, amplía el conjunto de transferencias posibles, altera el equilibrio financiero y a afecta los precios de los activos. Ahora bien, ¿a qué coste? Como cada consumidor propietario puede tener un perfil de riesgo distinto, entonces \textbf{no hay un objetivo único en la empresa} y maximizar el valor deja de tener sentido económico.
\end{la}

Por tanto, en dicho escenario nos sometemos a la imposibilidad de alcanzar un Equilibrio de RADNER con equivalencia ARROW-DEBREU (o, si existe, de una elevadísima complejidad analítica) o al vaciamiento de contenido de la empresa.

\subsection{Enfoque del Núcleo: Juegos cooperativos}






























3.1. El enfoque del núcleo e introducción de comportamiento estratégico
SCARF y DEBREU, apoyándose en el trabajo de EDGEWORTH [Ver Anexo IV], desarrollaron el enfoque del núcleo. 
Para ello, los autores pasan de la economía de EDGEWORTH formada por dos agentes que intercambian entre sí, a una economía formada por n agentes que pueden formar coaliciones entre ellos. Las coaliciones bloquearán el intercambio si al menos uno de sus miembros está mejor (y el resto es indiferente) intercambiando entre los miembros de la coalición que intercambiando con otra coalición. 
Lo que demuestran SCARF y DEBREU es que conforme va aumentando el número de miembros de la economía, aumenta tanto el tamaño de la coalición como las posibles coaliciones que pueden formarse. Esto lleva a que las asignaciones no bloqueadas por ninguna coalición sean cada vez menores, es decir, que el núcleo vaya siendo cada vez menor.
3.1.1. Modelo: SCARF y DEBREU
3.1.1.1. Supuestos
Supongamos que duplicamos la economía de intercambio de EDGEWORTH, de modo que:
	Pasan a existir dos agentes del tipo 1 y dos agentes del tipo 2. 
	Los agentes se han organizado en coaliciones. En concreto, hay una coalición formada por los dos consumidores de tipo 1 y por uno de los consumidores de tipo 2.
3.1.1.2. Desarrollo 
 
Los integrantes de la coalición estudian la posibilidad de intercambiar bienes entre sí a un precio relativo representado por la recta x̅C. 
Si, por ejemplo, los dos agentes del tipo 1 se sitúan en H entregarán 2FG del bien X_2 al agente de tipo 2 a cambio de recibir de este 2IJ unidades del bien X_1. Este último se situaría en el punto H^‘, sobre la recta de intercambio a una distancia de x̅ doble de la H y los dos agentes del tipo 1 en el punto H.
Esto significa que la coalición puede, con las dotaciones iniciales que tiene en su poder, y excluyendo del intercambio al otro agente, cambiar los bienes de forma tal que los dos consumidores del tipo 1 alcancen una curva de indiferencia U_1^1; y el de tipo 2 la U_1^2. Esto elimina al punto C, y como resulta obvio a otros muy próximos a él, como asignación perteneciente al núcleo de esta economía duplicada, pues implica menores niveles de utilidad para los agentes de la coalición que las posiciones H y H^‘ del Grafico anterior. Dicho de otro modo, la coalición bloquearía la asignación C porque podría, intercambiando entre sí sus dotaciones iniciales, mejorar la posición de todos sus miembros respecto a C.
Parece claro que a medida que aumenta el número de agentes en la economía, las asignaciones que pertenecen al núcleo de la misma se reducen. Así, SCARF y DEBREU demuestran entonces que si el número de individuos tiende a infinito, el núcleo irá reduciéndose progresivamente hasta alcanzar una única asignación. Esta será la misma que la que se obtenía bajo el equilibrio general walrasiano, dados los niveles de asignaciones iniciales.
3.1.1.3. Implicaciones
Conviene destacar que, al contrario de lo que podría parecer a simple vista, el enfoque de SCARF y DEBREU es más restrictivo que el de WALRAS, pues implica que dentro de las coaliciones los agentes con capaces de ponerse de acuerdo de forma perfecta e instantánea. No existe, por tanto, ningún tipo de coste de negociación,   a pesar de que podemos estar hablando de coaliciones de tamaños muy elevados. En otras palabras, se pasa de un equilibrio competitivo a otro cooperativo, similar a una economía asamblearia. 




3.4. Introducción del dinero
Dado que los modelos de equilibrio general sirven de nexo entre la microeconomía y la macroeconomía, resulta de interés introducir el dinero para poder estudiar cuestiones en materia monetaria. 
3.4.1. Modelo PATINKIN
En este sentido, destaca la obra de PATINKIN, quien se basó en la ecuación cuantitativa del dinero para introducir el mismo en los modelos de equilibrio general. 
3.4.1.1. Supuestos
En concreto, se parte de:
M·V=P·T
Siendo M la cantidad de dinero existente en la economía, V la velocidad de circulación del dinero, P el nivel de precios de la economía y T el volumen de transacciones que se realizan en la misma.
Así, los precios relativos de los bienes quedan determinados por los factores reales que intervienen en la economía, mientras que los precios absolutos dependen de factores monetarios. Para fijar la demanda de dinero de los agentes, se hace uso de la llamada Ecuación de Cambridge:
M=k·(∑_(n=1)^N▒p_n ·x_n )
Siendo k la inversa de la velocidad de circulación del dinero, tal que k=1/v
3.4.1.2. Desarrollo
A partir de esta ecuación es posible expresar la función de exceso de demanda del dinero de la siguiente forma:
z_(N+1) (p)=k·∑_(n=1)^N▒p_n ·x_n-M
La función z_(n+1) es homogénea de grado uno respecto al nivel de precios absolutos y la cantidad de dinero, por lo que, si todos los precios y la cantidad de dinero en la economía varían en la misma proporción, el exceso de demanda del dinero varía en la misma proporción.
No obstante, es problemático que dicha función sea homogénea de grado uno respecto a precios mientras que la función de exceso de demanda para el resto de los bienes es homogénea de grado cero, pues implicaría un comportamiento inconsistente de los consumidores:  Mantienen saldos monetarios para efectuar compras, por lo que el exceso de demanda de dinero es idénticamente igual al exceso de oferta de bienes, no tendría sentido entonces el diferente grado de homogeneidad. 
Para poder corregir esta contradicción, se introduce una ecuación adicional en el sistema que define el nivel de precios de una economía como la suma ponderada de los precios de los distintos bienes.
P=∑_(i=1)^I▒θ_i ·p_i
Donde θ representa la ponderaión de cada bien en la economía, de forma que ∑_i▒θ_i =1.
De esta forma, se llega a un sistema de n+2 ecuaciones:
	N funciones de exceso de demanda que toman la forma: Z_i=f_i (P_1/P,P_2/P,…,P_N/P,M/P)
	La función de exceso de demanda de dinero Z_(n+i)=f_(n+i) (P_1/P,P_2/P,…,P_N/P,M/P)-M/P  
	El nivel general de precios de la economía P=∑_(i=1)^I▒θ_i ·p_i
3.4.1.3. Implicaciones
Así:
	Todas las funciones de exceso de demanda son homogéneas de grado cero respecto a los precios absolutos y dinero;  al mismo tiempo, 
	Permite establecer que las funciones de demanda de bienes dependen de la cantidad de dinero existente en la economía y del nivel general de precios. 
En este sistema, un aumento de la cantidad de dinero provoca un aumento del nivel general de precios, pero los precios relativos no varían. 
Por tanto, a modo de conclusión, en este modelo ocurre que cambios en la cantidad de dinero circulando en la economía no tienen efectos reales. Estos cambios se limitan a cambiar el precio absoluto de los bienes en una proporción idéntica para todos ellos. 
La principal limitación del modelo es que la utilidad de los individuos no depende del dinero, por lo que no se realiza una modelización satisfactoria de la demanda del dinero. Si este no aporta utilidad a los agentes, ¿para qué quieren entonces demandarlo? 
Esta limitación propició el desarrollo de modelizaciones posteriores como los modelos MIU (Money in Utility) o los modelos CIA (cash in advance). 



3.5. Modelo de Equilibrio General Computable
Hasta la década de 1970, el análisis de equilibrio general era esencialmente teórico. Sin embargo, en 1969, el economista SCARF presentó un algoritmo numérico eficiente para el cálculo de equilibrios en economías complejas. Ello popularizó el uso de modelos de equilibrio general para conocer las características del equilibrio, así como para evaluar el impacto general de distintas opciones de política económica. A lo largo de los últimos años, se han producido dos avances que han permitido el rápido de desarrollo de los modelos de equilibrio general:
	Por un lado, los desarrollos posteriores del la Teoría del Equilibrio General que permiten expandirse para incluir varias características de los mercados en el mundo real. 
	En segundo lugar, el desarrollo de la capacidad de los ordenadores y de software electrónico cada vez más potentes ha hecho posible crear modelos enormemente complejos, que incluyen un número muy elevado de bienes y tipos de hogares.  


3.5.1. Metodología
La metodología empleada en los modelos de Equilibrio General Computable parte de:
	Fijación del número de bienes que se incluyen en el modelo; 
	Distintas formas posibles de las funciones de producción y de utilidad; , 
	Asignación de valores a determinados parámetros de las funciones en base a la evidencia empírica; y,
	Por último, un modelo de equilibrio general completo debe especificar cómo interviene el Estado en el mismo.  
En definitiva, el modelo debe especificar la totalidad de los flujos, tanto de fuentes como de usos de renta, que caracterizan la economía objeto de modelización. 
Más adelante, los resultados obtenidos son contrastados con la realidad. En función del resultado, el modelo es calibrado para hacerlo más preciso. Finalmente, una vez el modelo está listo, se lleva a cabo un análisis comparativo introduciendo distintos escenarios de política económica para ver los efectos de las mismas. Este análisis permitirá a las autoridades conocer de antemano el impacto de las distintas opciones de política económica, y en base a los resultados optar por una u otra opción. Una vez determinados la tecnología (oferta) y las preferencias (demanda), se deben obtener tanto los precios como cantidades que caracterizan el equilibrio general. Se considera que debe existir un equilibrio, pero éste no es siempre fácil de obtener, especialmente cuando el número de bienes y hogares incluidos en el modelo es elevado.
3.5.2. Resolución: algoritmo de SCARF
Los modelos de equilibrio general emplean ordenadores y software avanzando para la obtención de resultados. Para ello, se utiliza el algoritmo de SCARF con una serie de modificaciones posteriores, que busca el equilibrio de vaciado de mercado reproduciendo la forma en la que funcionan los mercados en la vida real.  
3.5.3. Aplicaciones
3.5.3.1. Comercio Internacional
Uno de los usos pioneros de este tipo de modelos fue el estudio del impacto de las barreras comerciales. Gran parte del debate sobre dichas barreras se ha centrado en su impacto sobre los salarios reales, por lo que los modelos de equilibrio general parecen especialmente apropiados para el estudio de dicho fenómeno.
Algunos ejemplos concretos son:
	Acuerdo de libre cambio de América del Norte (NAFTA);  
	Acuerdo Transatlántico para el Comercio y la Inversión (TTIP).  
3.5.3.2. Impuestos y transferencias
Los modelos del equilibrio general Computable también sirven para evaluar el impacto de un cambio en las políticas impositivas y de transferencias. 
Así, el efecto que un cambio en el impuesto sobre la renta puede tener sobre la oferta de trabajo se estudia a través de estos modelos. Cambios en el nivel de transferencias pueden afectar al nivel de ahorro e inversión por parte de los agentes, por lo que es necesario ser especialmente cuidadoso a la hora de elaborar el modelo (por ejemplo diferenciando a los consumidores por rangos de edad para evaluar los efectos de los distintos programas de jubilación).
3.5.3.3. Políticas medioambientales
Los detractores de las políticas de protección del Medio Ambiente siempre mencionan el supuestamente elevado coste económico y en términos de bienestar que implica la adopción de dichas medidas. Los modelos de equilibrio general permiten cuantificar el impacto real de estas medidas, y por lo tanto reducir la incertidumbre existente a este respecto. Así, especificando los objetivos medioambientales derivados de un menor uso de un determinado producto (por ejemplo la emisión de gases de efecto invernadero), es posible emplear estos modelos para estudiar el coste económico de distintas estrategias existentes para alcanzar el objetivo, eligiéndose la menos costosa.  
3.5.3.4. Políticas de planificación regional y urbanística 
También se utilizan para examinar cuestiones de política económica relacionadas con aspectos de localización espacial. La construcción de dichos modelos requiere una especial atención en variables como los costes de transporte de los bienes y los costes de movilidad del trabajo. , 
3.5.3.5. España
En España se utiliza modelo de equilibrio general dinámico y estocástico EREMS, elaborado por el Ministerio de Economía, la Universidad de Valencia, Fundación Rafael del Pino y BBVA Research. 
	Su funcionamiento está determinado por un conjunto de ecuaciones, variables y parámetros que describen cómo se comporta a nivel agregado la economía española.
La mayor parte de las variables son endógenas en el modelo. Cuando la economía recibe perturbaciones, como las originadas por el coronavirus, la respuesta de estas variables endógenas ofrece una imagen de la dinámica de los principales agregados macroeconómicos que suceden a estas perturbaciones.

 







\clearpage
\section{Conclusión} 

\subsection{Extensiones y relación con otras partes del Temario}


\subsection{Opinión}



\subsubsection{Idea final (Salida o cierra)}


\clearpage
\pagenumbering{gobble}
\MyBibliography
\MyBibItem{Autor}{Título}{Año}{DOI -- LINK}




%ANEXOS
\clearpage
\anexo{\textbf{Preguntas Test}}
%\input{\TemaCode/questions.tex} 


\clearpage
\anexo{Teoría del Equilibrio General (producción): ARROW-DEBREU}
\clearpage

La Teoría del Equilibrio General formulada por ARROW-DEBREU, con producción, conlleva hacer los siguientes cambios.

Para ello, parte de los siguientes supuestos. Relativos a los fundamentales de la economía, se se asume que:
\begin{lnum}
    \item \textbf{Productores}. Suponemos que cada productor viene caracterizada por una tecnología o conjunto de producción, $Y_j\subset\mathbb R^L$. Asumimos tambiém que los conjuntos de producción de toda industria son cerrados, acotados y estríctamente convexos. Es decir, dan lugar a Funciones de Producción Neoclásicas de Buen Comportamiento;\footnote{%
        Cuando los conjuntos de producción no son estrictamente convexos, las cosas se vuelven más complicadas porque las correspondencias $y_j(p)$ pueden dejar de ser unívocas. Una situación productiva de considerable importancia teórica y práctica, y que ciertamente no queremos descartar por hipótesis, es el caso de rendimientos constantes a escala. Con rendimientos constantes, sin embargo, los conjuntos de producción no son ni estrictamente convexos ni acotados superiormente (salvo en el caso trivial en el que no pueda producirse ninguna cantidad positiva de ningún bien). En principio, aún podríamos considerar los equilibrios como los ceros de una \textit{correspondencia de exceso de demanda inclusiva de la producción}, definida como en para un subconjunto de precios estrictamente positivos. Las correspondencias, no obstante, no constituyen buenos sistemas de ecuaciones (por ejemplo, no pueden diferenciarse). Por ello, en tales casos suele resultar mucho más conveniente caracterizar los equilibrios como las soluciones de un sistema ampliado de ecuaciones que involucra los lados de la producción y del consumo de la economía.
    }
    \item \textbf{Consumidores}. Cada consumidor viene caracterizado por un conjunto de consumo $X_i\subset\mathbb R^L$ y una relación de preferencias $\succsim_i$ definida dentro de este conjunto de consumo. Asumimos que estas preferencia son racionales (completas y transitivas);
    \item \textbf{Dotaciones iniciales}. Cada consumidor dispone, inicialmente, de una dotación o renta constituida por: (i) una dotación de recursos inicial $\omega_i\in \mathbb R^L$; y, (ii) las rentas obtenidas por su participación en los beneficios de cada empresa, $0\leq\theta_{ij}$ (teniendo en cuenta $\sum_{i=1}^I\theta_{ij}=1$). Nótese que nada impide la distribución extremadamente desigual de las participaciones empresariales, pues la restricción se impone únicamente sobre la suma de las participaciones para cada empresa, no para cada consumidor. 
\end{lnum}

Por lo que respecta a los supuestos institucionales, supondremos que en el equilibrio se fija un precio para cada bien, incluidos aquellos que no se comerciarán, y que no existen restricciones informativas.

Por último, desde una perspectiva conductual de los agentes, suponemos que los consumidores y empresas toman los precios como dados, son precio-aceptantes. 

Bajo este marco formal, podemos entonces definir el Equilibrio General Competitivo o Equilibrio General de mercado como la conjunción de las siguientes condiciones:

$$\begin{aligned}
    &\text{Dada una economía descentralizada, definida por:}\qquad \left(\{(X_i,\succsim_i)\}_{i=1}^I,\;\{Y_j\}_{j=1}^J,\;\{\omega_i,\theta_{i1},\dots \theta_{iJ}\}_{i=1}^I\right)\\[2em]
    &((\mathbf x^*,\mathbf y^*),\;\mathbf p=(p_1,p_2,\dots,p_L))\equiv \underset{\text{\scriptsize Competitivo o de mercado}}{\text{Equilibrio}}\iff\;
    \left\{\begin{aligned}
        &\text{(1)}&& \forall j,\quad\;y_j^*\in Y_j\;\;\text{maximiza $\Pi_j$},\;\;\;\text{i.e.},\\[0.5em]
        &&&\hspace{4em}\mathbf p\cdot y_j\leq \mathbf p\cdot y_j^*,\;\;\;\forall y_j\in Y_j\\[1em]
        &\text{(2)}&& \forall i,\quad\;x_i^*\in X_i\;\;\text{es máximo para $\;\succsim_i\;\;$ restringido a la dotación}\\[0.5em]
        &&&\hspace{4em}\left\{x_i\in X_i:\;\;\mathbf p\cdot x_i\leq \mathbf p\cdot \omega_i+\sum_{j=1}^J\theta_{ij}\;\mathbf p\cdot y_j^*\right\}\\[1em]
        &\text{(1)}&& \sum_{i=1}^Ix_i^*=\sum_{i=1}^I\omega_i+\sum_{j=1}^Jy_j^*
    \end{aligned}\right.
\end{aligned}$$


La primera condición establece que, en un equilibrio walrasiano, las empresas maximizan sus beneficios dados los precios de equilibrio $\mathbf p$. La segunda establece que los consumidores maximizan su bienestar dados, en primer lugar, los precios de equilibrio y, en segundo lugar, la riqueza derivada de sus dotaciones de bienes y de sus participaciones en los beneficios. Por último, la última condición establece que los mercados deben vaciarse en un equilibrio; esto es, todos los consumidores y las empresas deben poder realizar los intercambios deseados a los precios de mercado vigentes. Por tanto, la noción de equilibrio walrasiano en una economía de mercado es la predicción positiva del resultado de un sistema de mercados en los que consumidores y empresas son precio-aceptantes y la riqueza de los consumidores deriva de sus dotaciones iniciales de insumos y participación en los beneficios de la empresa. 

No obstante, a efectos del análisis formal, resulta sumamente útil poder expresar los equilibrios como las soluciones de un sistema de ecuaciones. Objetivo que perseguiremos en adelante, aspirando a ser muy concretos pero debiendo recurrir a la imposición de supuestos más fuertes para simplificar el análisis.

Para ello introducimos la noción de \textbf{función de exceso de demanda}. Definimos la Función de Exceso de Demanda del consumidor $i$ dado el vector de precios $\mathbf p$, como:

$$\begin{aligned}
    &\text{Sea,}\;\;\underbrace{m_i(\mathbf p)=\mathbf p\cdot \omega_i+\sum_{j=1}^J\theta_{ij}\Pi_j(\mathbf p)}_{\text{\scriptsize Riqueza/Ingreso walrasiano del consumidor: $m_i$}}\;\;\Longrightarrow\;\;\underset{\substack{\text{\scriptsize Vector de demanda neta/\textit{excesiva} del consumidor $i$ para cada bien,}\\\text{\scriptsize por encima/debajo de la dotación inicial de insumos productivos}}}{x_i(\mathbf p,\;m_i(\mathbf p))-\omega_i\in \mathbb R^L}\\[1em]
    &\text{FED Individual: }\;\underset{\text{\scriptsize \textit{Lo que finalmente tengo menos lo que inicialmente tenía}}}{\boxed{z_i(\mathbf p)\in\mathbb R^L\equiv x_i\left(\mathbf p,\;m_i(\mathbf p)\right)-\omega_i}}\\[1em]
    &\text{FED Agregada: }\;\underset{\substack{\\[0.5em]\\\text{\scriptsize El vector/cesta de consumo resultante menos el conjunto de dotaciones iniciales y el conjunto de bienes producidos,}\\\text{\scriptsize donde 
    $\left\{\begin{aligned}
        &z(\mathbf p)>0\sim \text{Exceso de demanda generalizado}\\
        &z(\mathbf p)\leq0 \sim \text{Plétora generalizada de recursos}
    \end{aligned}\right.$}}}{\boxed{z(\mathbf p)\in\mathbb R^L\equiv \sum_{i=1}^I\left[x_i\left(\mathbf p,\;m_i(\mathbf p)\right)-\omega_i\right]-\sum_{j=1}^Jy_j(\mathbf p)=\sum_{i=1}^Iz_i(\mathbf p)-\sum_{j=1}^Jy_j(\mathbf p)}}\\[2em]
    &\text{Entonces,}\qquad \boxed{\mathbf p\equiv \text{Equilibrio}\;\;\iff\;\;z(\mathbf p)\leq 0}
\end{aligned}$$

Por tanto, la demanda neta se define como lo que el consumidor quiere consumir menos lo que trae como dotación, ya que las ganancias entran en el presupuesto como ingreso, no como vector de bienes. Y, de manera general, la Función de Exceso de Demanda establece las condiciones del conjunto de bienes resultantes de la interacción, teniendo cuenta las dotaciones iniciales, las demandas individuales, y la tecnología de producción respecto a los bienes producidos. Además, por congruencia, al exigir la imposibilidad de una plétora general de bienes, es decir, de un exceso generalizado de bienes, se impone la condición de que, dado un vector de precios, el conjunto de bienes resultante genere un exceso de demanda agregado negativo o nulo. Esto es lo que se conoce como la condición de \textbf{factabilidad agregada de la economía}, $\sum_{i}x_i^*\leq \sum_{i}\omega_i+\sum_jy_j^*$: no puede haber demanda neta positiva agregada si no hay recursos suficientes.

No obstante, con preferencias localmente insaciables, si el precio de un bien es positivo ($p_\ell>0$), no puede sobrar en equilibrio, porque habría posibilidad de aumentar beneficios/valor o habría desequilibrio: $p_\ell\geq0,\;\;z_\ell(\mathbf p)\leq 0,\;\; p_\ell z_\ell(\mathbf p)=0$. Entonces, necesariamente cualquier bien $\ell$ puede presentar oferta excesiva ($z_\ell(\mathbf p)<0$) pero solo si este bien es gratuito: $p_\ell=0$.\footnote{%
    Como ejemplo, el bien $\ell$ puede ser un \textit{mal}. Entonces, podríamos esperar que su precio fuese cero, la demanda de los consumidores fuese cero y el exceso de oferta, dado por $z_\ell=\omega_\ell>0$, fuese eliminado mediante la condición de eliminación gratuita.
} Si este bien tiene un precio positivo, por el contrario, necesariamente debe vaciarse. 

De hecho, para simplificar aún más el desarrollo, si asumimos preferencias estríctamene monótonas, cualquier equilibrio general debe implicar un vector de precios estríctamente positivo, $\mathbf p\gg0$, ya que, de lo contrario los consumidores demandarían una cantidad infinita de cualquier bien gratuito. Por tanto, obtendríamos:

$$\begin{aligned}
    &\text{Si $\succsim_i$ son racionales, transitivas y estríctamente monótonas:}\\[1em]
    &\Longrightarrow\qquad \boxed{\mathbf p=(p_1,\dots p_L)\equiv \text{EGC}\iff z_\ell(\mathbf p)=0,\;\;\forall\ell=1,\dots, L:\quad z(\mathbf p)=0}\\[1em]
    &\text{Y, $z(\mathbf p)$ definida para todo vector $\mathbf p\gg 0$ cumple:}\\[0.5em]
    &\left\{\begin{aligned}
        &\text{(1)}&&z(\cdot )\in\text{Continua}\\
        &\text{(2)}&&z(\cdot )\in\mathcal H(0)\\
        &\text{(1)}&&\mathbf p\cdot z(\mathbf p)=0,\;\;\forall \mathbf p\\
        &\text{(1)}&&\exists s>0:z_\ell(\mathbf p)>-s,\;\;\forall \ell\;\&\;\forall\mathbf p\\
        &\text{(1)}&&\text{Si, }\;\;\mathbf p^n\to \mathbf p,\;\;\text{donde}\;\;\mathbf p\neq0\;\;\&\;\;p_\ell=0\;\;\text{para algún $\ell$},\\
        &&&\hspace{4em}\max\{z_1(\mathbf p^n),\dots ,z_L(\mathbf p^n)\}\to \infty
    \end{aligned}\right.
\end{aligned}$$

Con la excepción de la propiedad (v), el resto de propiedades son consecuencia directa de la definición y las propiedades paralelas de las funciones de demanda. La cota o límite $s$ se deriva de la no-negatividad de la demanda, que implica que la oferta neta de un consumidor al mercado de cualquier bien $\ell$ no puede ser mayor que su dotación inicial. La intuición detrás de la quinta propiedad es precisamente que cuando algún úprecio tiende a cero, un consumidor cuyo bienestar/utilidad tiende estríctamente a un límite positivo y cuyas preferencias son estríctamente monótonas, demandará una cantidad cada vez mayor de algún/os bien/es/commodity cuyo/s precio/s tiendan a cero, pero no de todos debido precisamente a la importancia de los precios relativos. 

De hecho, podems observar que, debido a las propiedades de la FDE agregada definida sobre el conjunto posible de vectores de precios estríctamente positivos, para verificar si un vector de precios estríctamente superior a cero es un vector de equilibrio (EGC), será suficiente con ver si es un vector de vaciamiento de mercados para todos salvo para uno:

$$\begin{aligned}
    &\text{Si $\mathbf p\gg 0$ es un vector de EGC, podemos comprobarlo:}\\[0.5em]
    &\Longrightarrow \mathbf p^*\gg 0:z_1(\mathbf p^*)=\cdots=z_{L-1}(\mathbf p^*)=0\\[0.5em]
    &\underset{\text{\scriptsize como,}}{\Longrightarrow}\;\;\mathbf p^*\cdot z(\mathbf p^*)=\sum_\ell p_\ell^*\cdot z_\ell(\mathbf p^*)=0\;\;\text{y,}\;\; p_L>0\\[0.5em]
    &\Longrightarrow z_L(\mathbf p^*)=0\underset{\substack{\text{\scriptsize denotamos vector FDE}\\\text{\scriptsize para todos los bienes desde 1}\\\text{\scriptsize hasta $L-1$, como}}}{\Longrightarrow}\hat z(\mathbf p^*)=(z_1(\mathbf p^*),\dots,z_{L-1}(\mathbf p^*))\\[0.5em]
    &\Longrightarrow \boxed{\mathbf{p}^*\gg 0\equiv \text{EGC}\iff \hat z(\mathbf p^*)=0}
\end{aligned}$$

Es decir, como la suma de los excesos de demanda debe ser cero, los mercados necesariamente son interdependientes. Por tanto, si $L-1$ mercados están en equilibrio, el último también lo estará. De manera análoga, si $L-1$ mercados presentan exceso de oferta, el último mercado necesariamente debe presentar exceso de demanda. 

De manera más general, este resultado nos permite obtener el precio de un bien en términos de los demás bienes. Es decir, permite determinar los precios de los bienes en términos relativos sin que se conozcan los valores absolutos de los precios. En conclusión, el equilibrio general competitivo es un conjunto de precios no negativos, asociado a un conjunto de demandas y ofertas de economías domésticas y empresas, tales que cada una de las demandas y ofertas es óptima y los excesos de demanda existentes son nulos (es decir, los planes de los agentes son compatibles y pueden llevarse a cabo).





\clearpage
\anexo{Teoría del Equilibrio General: Modelo de comercio secuancial (RADNER)}\label{anx:radner}

En el equilibrio básico de Radner, el “bien contingente” no es simplemente el numerario físico que existe en t=1; es un contrato en t=0 que promete entregar una unidad del numerario en un estado concreto. Aunque el numerario tenga precio 1 en todos los estados en t=1, no es el mismo bien económico recibir una unidad en un estado que en otro. Lo que se comercia en t=0 no es “el numerario hoy”, sino derechos contingentes a numerario mañana, estado por estado.

La transferencia de riqueza ocurre así: cuando un consumidor compra en t=0 una unidad del bien contingente del estado s, está pagando hoy para asegurarse más poder adquisitivo si ocurre s; cuando vende ese bien en corto, está haciendo lo contrario, renunciando a riqueza en s a cambio de más riqueza hoy o en otros estados. Por tanto, mediante posiciones largas o cortas en esos contratos, el consumidor reconfigura cómo se distribuye su riqueza entre los distintos estados futuros, sin cambiar el total “esperado” en un sentido físico, pero sí su perfil de riesgo.

El hecho de que el precio spot del numerario sea 1 en todos los estados no elimina esta función. Esa normalización solo significa que medimos toda la riqueza futura en unidades del numerario, es decir, en términos relativos a él. No implica que el numerario “no importe”, sino que lo usamos como unidad de cuenta. El precio relevante para la transferencia de riqueza no es el precio spot (que hemos fijado a 1), sino el precio a plazo q_s del contrato que paga numerario si ocurre el estado s. Esos q_s son los que reflejan cuánto vale hoy una unidad de riqueza en cada estado.

Respecto a si el numerario se comercia efectivamente en t=0: no se comercia como bien físico, pero sí se comercia como promesa contingente. Y no es necesario tener numerario como dotación inicial para participar: puedes vender el contrato en corto. Si haces eso, te comprometes a entregar numerario en t=1 si ocurre el estado correspondiente, financiándolo entonces con tus dotaciones o con comercio spot. Esa posibilidad de vender en corto es lo que hace posible la transferencia de riqueza incluso sin dotaciones iniciales del numerario.

En resumen: el bien contingente sirve como vehículo de transferencia de riqueza porque permite mover poder adquisitivo entre estados futuros; fijar su precio spot a 1 solo establece la unidad de cuenta; y el comercio relevante ocurre en (t=0\ a través de los contratos contingentes, no a través del numerario físico presente hoy.

Para ser lo más simples posible consideramos únicamente economías de intercambio. Además, tomamos $X_i=\mathbb{R}^L_+$ para todo $i$. Para comenzar, suponemos que hay dos fechas, $t=0$ y $t=1$, que no hay información alguna en $t=0$, y que la incertidumbre se resuelve completamente en $t=1$. 

\textsc{Ausencia de comercio ex-post: ¿Por qué esperar?}

Supongamos que en $t=0$ se establecen mercados para los $LS$ posibles bienes contingentes, y que $(x^1,\ldots,x^I)\in\mathbb{R}^{LSI}$ es una asignación de equilibrio de ARROW-DEBREU con precios $(p_{11},\ldots,p_{LS})\in\mathbb{R}^{LS}$. (Estos mercados son para la entrega de bienes en $t=1$, comúnmente se denominan mercados a plazo). Entonces, cuando llega el período $t=1$ se revela un estado del mundo $s$, los contratos se ejecutan y cada consumidor $i$ recibe $x^s_i=(x^s_{1i},\ldots,x^s_{Li})\in\mathbb{R}^L$. 

Imaginemos ahora que, después de esto pero antes del consumo efectivo de $x^s_i$, se abren en $t=1$ mercados para los $L$ bienes físicos (estos se llaman mercados al contado). ¿Habría algún incentivo para comerciar en estos mercados? La respuesta es \textbf{no}. Para ver por qué, supongamos que existieran ganancias potenciales del comercio entre los consumidores. Es decir, que existieran $X^s_i=(X^s_{1i},\ldots,X^s_{Li})$ para $i=1,\ldots,I$, tales que $\sum_i X^s_i\le \sum_i \omega^s_i$ y $(X^s_1,\ldots,X^s_I)\succeq (x^s_1,\ldots,x^s_I)$ para todo $i$, con al menos una preferencia estricta. Por la definición de optimalidad de Pareto, la asignación de equilibrio de Arrow-Debreu $(x^1,\ldots,x^I)\in\mathbb{R}^{LSI}$ no es óptima en el sentido de Pareto, contradiciendo la conclusión del primer teorema del bienestar. 

En resumen, en $t=0$ los consumidores pueden comerciar directamente hasta alcanzar una asignación globalmente óptima en el sentido de Pareto; por lo tanto, no hay razón para que tenga lugar comercio adicional. En otras palabras, la optimalidad de Pareto ex ante implica optimalidad de Pareto ex post y, por lo tanto, ausencia de comercio ex post.


\textsc{Mejora Paretiana: ¿Cómo \textit{obligar} el comercio ex post?}

Las cosas son diferentes si no todos los $LS$ mercados de bienes contingentes están disponibles en $t=0$. Entonces, el comercio inicial hacia una asignación óptima de Pareto puede no ser factible y es muy posible que ex post (es decir, después de la revelación del estado $s$) la asignación de consumo resultante no sea óptima en el sentido de Pareto. En ese caso, habría un incentivo para reabrir los mercados y volver a comerciar. Una posibilidad muy interesante, observada por primera vez por ARROW (1953), es que, incluso si no todos los bienes contingentes están disponibles en $t=0$, puede ocurrir aun así, bajo ciertas condiciones, que las posibilidades de recomercio en $t=1$ garanticen que se alcance la optimalidad de Pareto. Es decir, la posibilidad de comercio ex post puede compensar la ausencia de algunos mercados ex ante. 

Verificaremos que este es el caso siempre que al menos un bien físico pueda comerciarse de manera contingente en $t=0$ si, además, existen mercados al contado en $t=1$ y los precios de equilibrio al contado son anticipados correctamente en $t=0$. La intuición de este resultado es razonablemente directa: si el comercio al contado puede tener lugar dentro de cada estado, entonces la única tarea restante en $t=0$ es transferir eficientemente el poder adquisitivo total del consumidor entre estados. Esto puede lograrse utilizando comercio contingente en un solo bien. Mediante tal procedimiento somos capaces de reducir el número de mercados a plazo requeridos de $LS$ a $S$.

\textsc{Equilibrio de RADNER}

Seamos más específicos. En $t=0$ los consumidores tienen expectativas respecto de los precios al contado que prevalecerán en $t=1$ para cada estado posible $s\in S$. Denótese por $\mathbf p^s\in\mathbb{R}^L$ el vector de precios que se espera que prevalezca en el mercado al contado del estado $s$, y por el vector de expectativas global $\mathbf p=(p^1,\ldots,p^S)\in\mathbb{R}^{LS}$. Supóngase además que en la fecha $t=0$ existe comercio en los $S$ bienes contingentes denotados por $1_1$ a $1_S$; es decir, que existe comercio contingente únicamente en el bien físico con la etiqueta $1$. Denotamos por $q=(q_1,\ldots,q_S)\in\mathbb{R}^S$ el vector de precios de estos bienes contingentes comerciados en $t=0$.

Enfrentado a precios $q\in\mathbb{R}^S$ en $t=0$ y a precios al contado esperados $(p^1,\ldots,p^S)\in\mathbb{R}^{LS}$ en $t=1$, cada consumidor $i$ formula un plan de consumo, o de comercio, $(z^i_1,\ldots,z^i_S)\in\mathbb{R}^S$ para los bienes contingentes en $t=0$, así como un conjunto de planes de consumo en los mercados al contado $(x^i_1,\ldots,x^i_S)\in\mathbb{R}^{LS}$ para los diferentes estados que pueden ocurrir en $t=1$. Por supuesto, estos planes deben satisfacer una restricción presupuestaria. Sea $U_i(\cdot)$ una función de utilidad para $i$. Entonces, el problema del consumidor $i$ puede expresarse formalmente como

$$\begin{aligned}
    &\max_{\begin{aligned}
        (x_{1i,\ldots,x_{Si}})\in\mathbb R^{LS}_+\\[0.5em]
        (x_{1i},\ldots,x_{Si})\in\mathbb R^{S}
    \end{aligned}}\qquad Ui(x_{1i},\ldots,x_{Si})\\[1em]
    &\text{s.a.}\quad \left\{\begin{aligned}
        &\sum_{s}q_sz_{si}\leq0,\\
        &\mathbf p_s\cdot x_{si}\leq \mathbf p_s\cdot \omega_{si}+\mathbf p_{1s}z_{si},\quad \forall s
    \end{aligned}\right.
\end{aligned}$$

La restricción (i) es la restricción presupuestaria correspondiente al comercio en $t=0$. La familia de restricciones (ii) son las restricciones presupuestarias para los diferentes mercados al contado. Obsérvese que el valor de la riqueza en un estado $s$ está compuesto de dos partes: el valor de mercado de las dotaciones iniciales, $\mathbf p^s\cdot \omega^i_s$, y el valor de mercado de las cantidades $z^i_s$ del bien 1 compradas o vendidas a plazo en $t=0$. Obsérvese que no estamos imponiendo ninguna restricción sobre el signo o la magnitud de $z^i_s$. Si $z^i_s<-\omega^i_{1s}$, entonces se dice que en $t=0$ el consumidor $i$ está vendiendo el bien 1 en corto. Esto se debe a que está vendiendo en $t=0$, contingente a que ocurra el estado $s$, más de lo que tiene en $t=1$ si $s$ ocurre. Por lo tanto, si $s$ ocurre, tendrá que comprar en el mercado al contado la cantidad adicional del primer bien requerida para el cumplimiento de sus compromisos. La posibilidad de vender en corto está, sin embargo, limitada indirectamente por el hecho de que el consumo, y por lo tanto la riqueza ex post, debe ser no negativa para todo $s$.

Para definir una noción apropiada de comercio secuencial impondremos una condición clave: las expectativas de los consumidores deben ser autocumplidas, o racionales; es decir, requerimos que las expectativas de los consumidores acerca de los precios que despejarán los mercados al contado para los diferentes estados $s$ efectivamente los despejen una vez que llega la fecha $t=1$ y se revela un estado $s$.

\textbf{Definición.-} (Equilibrio de RADNER) Una colección formada por un vector de precios $q=(q_1,\ldots,q_S)\in\mathbb{R}^S$ para bienes contingentes del primer bien en $t=0$, un vector de precios al contado para cada $s$, y, para cada consumidor $i$, planes de consumo $z_i=(z^i_1,\ldots,z^i_S)\in\mathbb{R}^S$ en $t=0$ y $x_i=(x^i_1,\ldots,x^i_S)\in\mathbb{R}^{LS}$ en $t=1$, constituye un equilibrio de RADNER si:
\begin{lnum}
    \item Para todo $i$, los planes de consumo $z_i,x_i$ resuelven el problema; y,
    \item $\sum_iz_{si}^*\leq 0$ y $\sum_{i}x_{si}^*\leq\sum_i\omega_{si}$ para todo $s$
\end{lnum}

En un equilibrio de RADNER, el comercio tiene lugar a lo largo del tiempo y, en contraste con el entorno de ARROW-DEBREU, los agentes económicos enfrentan una secuencia de conjuntos presupuestarios, uno en cada fecha-estado.

Podemos ver que todas las restricciones presupuestarias son homogéneas de grado cero con respecto a los precios. Esto significa que los conjuntos presupuestarios permanecen inalterados si el precio de un bien físico en cada fecha-estado (es decir, un precio para cada conjunto presupuestario) se normaliza arbitrariamente a $1$. Es natural elegir el primer bien y fijar $p^s_1=1$ para todo $s$, de modo que una unidad del bien contingente del estado $s$ paga entonces 1 dólar en el estado $s$. Nótese que esto aún deja un grado de libertad, el correspondiente a los intercambios a plazo en la fecha $0$ (por lo que podríamos fijar $q_1=1$, o quizás $\sum_s q_s=1$).

El resultado clave de este desarrollo (que no se muestra pero puede verse en [\authorfont{Ver Manual}]) es que: el conjunto de asignaciones de equilibrio de ARROW-DEBREU (inducidas por la organización de comercio de una sola vez en $LS$ bienes contingentes) y el conjunto de asignaciones de equilibrio de RADNER (inducidas por comercio contingente en un solo bien, seguido secuencialmente por comercio al contado) son idénticos.

Es importante enfatizar que, aunque el concepto de equilibrio de Radner reduce el número de bienes contingentes requeridos para alcanzar la optimalidad (de $LS$ a $S$), esta reducción no se obtiene de manera gratuita. Con el menor número de contratos a plazo, la anticipación correcta de los precios futuros al contado se vuelve crucial.

\clearpage
\anexo{Activos: Definición y ejemplos}\label{anx:activos}
\textbf{Definición.-} Una unidad de un activo, o valor, es un título para recibir una cantidad $r_s$ del bien $1$ en la fecha $t$ si ocurre el estado $s$. Por lo tanto, un activo se caracteriza por su vector de rendimientos $r=(r_1,\ldots,r_S)\in\mathbb{R}^S$.

Si los rendimientos son en bienes físicos, el activo se denomina real (una pieza duradera de maquinaria o un contrato de futuros para la entrega de cobre serían ejemplos). Si son en dinero fiduciario, se denominan financieros (un bono del gobierno, por ejemplo). También son posibles casos mixtos. Ejemplos de activos incluyen los siguientes:
\begin{lnum}
    \item $r=(1,\ldots,1)$. Este activo promete la entrega futura no contingente de una unidad del bien $1$. Sus contrapartidas en el mundo real son los mercados de futuros de materias primas. En el caso especial en el que existe un único bien de consumo (es decir, $L=1$), llamamos a este activo el activo seguro (o libre de riesgo). Es importante darse cuenta de que, con más de un bien físico, un contrato de futuros no es libre de riesgo: su rendimiento en términos de poder adquisitivo depende de los precios al contado/spot de todos los bienes;
    \item $r=(0,\ldots,0,1,0,\ldots,0)$. Este activo paga una unidad del bien $1$ si y solo si ocurre un determinado estado. En el entorno teórico actual se los denomina a menudo \textit{activos/valores de Arrow}; y,
    \item $r=(1,2,1,2,\ldots,1,2)$. Este activo paga una unidad incondicionalmente y, además, otra unidad en los estados etiquetados como pares.
\end{lnum}

Una opción también podría considerarse un ejemplo. Este es activo derivado, es decir, un activo cuyos rendimientos se derivan de los rendimientos de otro activo (por ejemplo, $r$). Supóngase que existe un activo primario con vector de rendimientos $r\in\mathbb{R}^S$. Entonces, una opción de compra sobre el activo primario con precio de ejercicio $c\in\mathbb{R}$ es ella misma un activo. Una unidad de este activo otorga la opción de comprar, después de que se revela el estado (pero antes de que se paguen los rendimientos), una unidad del activo primario al precio $c$ (el precio $c$ está expresado en unidades del numerario, es decir, del bien $1$). ¿Cuál es el vector de rendimientos $r(c)$ de la opción? En un estado dado $s$, la opción se ejercerá (racionalmente) si y solo si $r_s>c$ (descuidamos el caso $r_s=c$). Por lo tanto, $r(c)=(\max\{0,r_1-c\},\ldots,\max\{0,r_S-c\})$. Por ejemplo, un activo primario con rendimientos $r=(4,3,2,1)$, algunos ejemplos serían: $r(3.5)=(0.5,0,0,0)$, $r(2.5)=(1.5,0.5,0,0)$, $r(1.5)=(2.5,1.5,0.5,0)$.


\end{document}